\documentclass{article}  % if you need to pass options to natbib, use, e.g.: %     \PassOptionsToPackage{numbers, compress}{natbib} % before loading neurips_2026  % The authors should use one of these tracks. % Before accepting by the NeurIPS conference, select one of the options below. % 0. "default" for submission \usepackage{neurips_2026} % the "default" option is equal to the "main" option, which is used for the Main Track with double-blind reviewing. % 1. "main" option is used for the Main Track %  \usepackage[main]{neurips_2026} % 2. "position" option is used for the Position Paper Track %  \usepackage[position]{neurips_2026} % 3. "eandd" option is used for the Evaluations & Datasets Track  % \usepackage[eandd]{neurips_2026}  % if you need to opt-in for a single-blind submission in the E&D track:  %\usepackage[eandd, nonanonymous]{neurips_2026} % 4. "creativeai" option is used for the Creative AI Track %  \usepackage[creativeai]{neurips_2026} % 5. "sglblindworkshop" option is used for the Workshop with single-blind reviewing  % \usepackage[sglblindworkshop]{neurips_2026} % 6. "dblblindworkshop" option is used for the Workshop with double-blind reviewing %  \usepackage[dblblindworkshop]{neurips_2026}  % After being accepted, the authors should add "final" behind the track to compile a camera-ready version. % 1. Main Track  % \usepackage[main, final]{neurips_2026} % 2. Position Paper Track %  \usepackage[position, final]{neurips_2026} % 3. Evaluations & Datasets Track  % \usepackage[eandd, final]{neurips_2026} % 4. Creative AI Track %  \usepackage[creativeai, final]{neurips_2026} % 5. Workshop with single-blind reviewing %  \usepackage[sglblindworkshop, final]{neurips_2026} % 6. Workshop with double-blind reviewing %  \usepackage[dblblindworkshop, final]{neurips_2026} % Note. For the workshop paper template, both \title{} and \workshoptitle{} are required, with the former indicating the paper title shown in the title and the latter indicating the workshop title displayed in the footnote. % For workshops (5., 6.), the authors should add the name of the workshop, "\workshoptitle" command is used to set the workshop title. % \workshoptitle{WORKSHOP TITLE}  % "preprint" option is used for arXiv or other preprint submissions  % \usepackage[preprint]{neurips_2026}  % to avoid loading the natbib package, add option nonatbib: %    \usepackage[nonatbib]{neurips_2026}  \usepackage[utf8]{inputenc} % allow utf-8 input \usepackage[T1]{fontenc}    % use 8-bit T1 fonts \usepackage{hyperref}       % hyperlinks \usepackage{url}            % simple URL typesetting \usepackage{booktabs}       % professional-quality tables \usepackage{amsfonts}       % blackboard math symbols \usepackage{nicefrac}       % compact symbols for 1/2, etc. \usepackage{microtype}      % microtypography \usepackage{xcolor}         % colors  \usepackage{amsmath} \usepackage{amssymb} \usepackage{mathtools} \usepackage{amsthm} \usepackage{multirow}  \usepackage{hyperref} \usepackage{url} \usepackage{float}  \usepackage{array} \usepackage{microtype} \usepackage{graphicx} \usepackage{xcolor} % \usepackage{subfigure} \usepackage{amsmath} % \usepackage{breqn} \usepackage{bm} \usepackage{dsfont} \usepackage{mathtools} \usepackage{nccmath} \usepackage{amssymb} \usepackage{multirow} \usepackage{booktabs} % for professional tables \usepackage{algorithm} \usepackage{algpseudocode} \usepackage{varwidth} \usepackage{pifont} \usepackage{minted} \usepackage{makecell} \usepackage[font=small,labelfont=bf,tableposition=top]{caption} \usepackage{wrapfig} \usepackage[flushleft]{threeparttable} % \usepackage[vlined,linesnumbered,ruled,algo2e]{algorithm2e} % \usepackage[linesnumbered,ruled,vlined]{algorithm2e} % \newcommand\mycommfont[1]{\footnotesize\ttfamily\textcolor{blue}{#1}} % \SetCommentSty{mycommfont} \usepackage{subcaption} \usepackage{varwidth} \usepackage{pifont} \usepackage{makecell} \usepackage{enumitem} \usepackage{soul} \usepackage[normalem]{ulem} % For underlining with \ul \usepackage{listings} % Adjust the underline height with \setul \setul{0.3ex}{0.6pt} % First argument: distance from text; second: underline thickness % \usepackage{footnote} % \makesavenoteenv{tabular}  % if you use cleveref.. \usepackage[capitalize,noabbrev]{cleveref} % --- Preamble --- \usepackage{tcolorbox} \usepackage{inconsolata} \tcbuselibrary{breakable, skins}  \newtcolorbox{promptbox}[1][]{   breakable,   colback=gray!8,   colframe=gray!50,   fonttitle=\bfseries\small,   title={Prompt},   #1 }  % --- Preamble --- \usepackage{tcolorbox} \usepackage{listings} \usepackage{xcolor} \tcbuselibrary{breakable, skins}  \lstdefinestyle{cstyle}{   language=C,   basicstyle=\ttfamily\scriptsize,   keywordstyle=\color{blue!70!black}\bfseries,   commentstyle=\color{green!50!black}\itshape,   stringstyle=\color{orange!80!black},   numbers=left,   numberstyle=\tiny\color{gray},   numbersep=8pt,   breaklines=true,   tabsize=4,   showstringspaces=false,   captionpos=b, }  \newtcolorbox{codebox}[1][]{   breakable,   colback=gray!5,   colframe=gray!50,   fonttitle=\bfseries\small,   title={Code},   #1 } %%%%%%%%%%%%%%%%%%%%%%%%%%%%%%%% % THEOREMS %%%%%%%%%%%%%%%%%%%%%%%%%%%%%%%% \theoremstyle{plain} \newtheorem{theorem}{Theorem}[section] \newtheorem{proposition}[theorem]{Proposition} \newtheorem{lemma}[theorem]{Lemma} \newtheorem{corollary}[theorem]{Corollary} \theoremstyle{definition} \newtheorem{definition}[theorem]{Definition} \newtheorem{assumption}[theorem]{Assumption} \theoremstyle{remark} \newtheorem{remark}[theorem]{Remark}    \newcommand{\wuyang}[1]{{\color{purple}\sf{[Wuyang: #1]}}}   % Note. For the workshop paper template, both \title{} and \workshoptitle{} are required, with the former indicating the paper title shown in the title and the latter indicating the workshop title displayed in the footnote.  % \title{KissatEvolve: LLM-based Evolutionary Technique for Evolving Large-Scale SAT Solver with\\Densely Annotated Memory Bank} % \title{KissatEvolve: Controllable and Scalable Synthesis of SAT Solver Heuristics with Densely Annotated Memory Bank} \title{KissatEvolve: Controllable and Scalable Synthesis of SAT Heuristics with Densely Annotated Memory Bank}   % The \author macro works with any number of authors. There are two commands % used to separate the names and addresses of multiple authors: \And and \AND. % % Using \And between authors leaves it to LaTeX to determine where to break the % lines. Using \AND forces a line break at that point. So, if LaTeX puts 3 of 4 % authors names on the first line, and the last on the second line, try using % \AND instead of \And before the third author name.   \author{%   David S.~Hippocampus\thanks{Use footnote for providing further information     about author (webpage, alternative address)---\emph{not} for acknowledging     funding agencies.} \\   Department of Computer Science\\   Cranberry-Lemon University\\   Pittsburgh, PA 15213 \\   \texttt{hippo@cs.cranberry-lemon.edu} \\   % examples of more authors   % \And   % Coauthor \\   % Affiliation \\   % Address \\   % \texttt{email} \\   % \AND   % Coauthor \\   % Affiliation \\   % Address \\   % \texttt{email} \\   % \And   % Coauthor \\   % Affiliation \\   % Address \\   % \texttt{email} \\   % \And   % Coauthor \\   % Affiliation \\   % Address \\   % \texttt{email} \\ }   \begin{document}   \maketitle   \begin{abstract} % We present KissatEvolve, an LLM-based evolutionary search method, with a novel densely annotated memory bank, aimed at totally automatic, controllable, and scalable synthesis of SAT solver heuristics. % Modern CDCL SAT solvers, such as Kissat\cite{biere2025cadical}, employ a large number of heuristics whose configuration space is prohibitively large to search exhaustively. % Existing approaches that leverage large language models (LLMs) for solver design have only explored a narrow slice of this space, discovering a limited set of solvers or heuristic variants, and failing to scale up the coverage of the design space. % Furthermore, these approaches overlook the value of producing specialized solvers tailored to specific sub-classes of benchmark problems.  % Our KissatEvolve framework is the first to bring both scalability and controllability to LLM-generated SAT heuristics. % %  % First, to support autonomous exploration of a large and diverse space of heuristic variants, KissatEvolve combines a self-adaptive agent workflow with a highly parallelized generation-and-evaluation pipeline. % The result is a fully automated system that requires minimal human intervention while substantially reducing the wall-clock time needed to explore the heuristic design space. % %  % Second, KissatEvolve maintains a memory bank that tracks both high- and low-performing generated heuristics, % each annotated with rich semantics, including solver analyses, search statistics, and traces of code diffs. % This memory bank not only improves token-level sampling efficiency, but also conditions KissatEvolve on selectively retrieved exemplars, enabling fine-grained control over solver behavior in subsequent heuristic generations. % %  % {\color{red} % Using KissatEvolve, we conducted a comprehensive study spanning 25 functions across 6 SAT solvers, % from which we discovered 5 promising heuristics, % improving the PAR-2 scores by XXX\% over state-of-the-art baselines, % and, when controllability is enabled, achieving an additional XXX\% PAR-2 improvement on targeted SAT, UNSAT, easy, and hard instance subsets. % } %%%%%%%%%%%%%%%% We present \textbf{KissatEvolve}, an LLM-based framework for fully automatic, scalable, and controllable synthesis of SAT solver heuristics. %  Modern CDCL solvers such as Kissat employ a large ensemble of tightly coupled heuristic components whose joint configuration space is prohibitively large to search exhaustively. Existing LLM-based approaches explore only a narrow slice of this space, confining search to a single solver or a small set of heuristic functions, and optimizing exclusively for aggregate PAR-2 without regard for targeted problem subcategories. %  KissatEvolve is the first framework to address both limitations simultaneously. For \textbf{scalability}, KissatEvolve combines a self-adaptive agentic workflow with a fully parallelized generation-and-evaluation pipeline that spans multiple solver families and heuristic functions concurrently, substantially reducing the wall-clock time required to cover the design space while minimizing human intervention. For \textbf{controllability}, KissatEvolve maintains a densely annotated memory bank that records both high- and low-performing heuristic variants, each enriched with solver analyses, search statistics, and code-diff traces. Selectively retrieved exemplars condition subsequent LLM generations toward user-specified subcategories (SAT, UNSAT, easy, or hard problems), enabling fine-grained steering of solver behavior.  Using KissatEvolve, we conduct a comprehensive study of 25 heuristic functions across 6 SAT solvers, finding improved heuristics for 3 solvers that reduce PAR-2 by 12.80\% on average over state-of-the-art baselines and by up to 43.76\% on targeted instance subcategories with controllability enabled. We provide code and generated solvers in the supplementary material, and promise to release them to the public upon acceptance. \end{abstract}   % stronger justification: why the specific/detailed design of memory bank helpful — help % name1: wrap up framework (techniques) % name2: solver   \section{Introduction}  % In modern CDCL SAT solvers, overall performance is strongly affected by the design of heuristic components such as restarting, score rescaling, variable activity updates, phase selection, and conflict analysis. % These components interact in complex ways, and effective improvements usually require repeated manual tuning and substantial solver expertise. % %  % Our goal is to replace manual engineering efforts with a large-scale evolution algorithm combined with LLM-based strategy generation.  % Traditional evolutionary agent frameworks lack structured memory, and therefore often revisit similar low-value ideas; % without explicit mutation control, many generated variants also differ only in vague or overlapping ways. % As the candidate population grows, orchestration, solver building, and evaluation further become serious system bottlenecks. % Our system is designed for this demanding setting of large-scale parallel evolution across many solver families and target functions, with explicit reuse of historical experience.   % To address these fundamental bottlenecks, we investigate two central questions: % \vspace{-0.5em} % \begin{center} % \fbox{ %     \parbox{0.9\linewidth}{ %         \textit{\textbf{Q1:} Can we explore better SAT heuristics by evolving \textcolor{red}{large-scale} SAT solvers?} \\ %         \textit{\textbf{Q2:} Can we make LLM exploration controllable, such that we can generate heuristics specifically tailored to selected problems?} %     } % } % \end{center}    % To address these challenges, we contribute in three directions: % controlled mutations for more precise local exploration, % a scalable asynchronous pipeline that increases the number of SAT solvers evolved in parallel, % and a densely annotated memory bank that converts previous generations into reusable optimization experience. % Together, these additions transform the framework from a single-generation search loop into a scalable, retrieval-augmented evolutionary system for SAT solver improvement.    Modern conflict-driven clause learning (CDCL) SAT solvers, exemplified by Kissat~\cite{biere2025cadical} and its derivatives, owe much of their performance to a tightly coupled ensemble of heuristic components: restart schedules, variable activity decay and rescaling, branching and phase selection, clause learning and deletion policies, and conflict analysis variants, among others. These components interact in subtle, problem-dependent ways, and advancing additional performance on top of state-of-the-art solvers typically demands expert intuition, careful manual tuning, and extensive empirical validation across diverse benchmark families.  % Historically, the cumulative effect of many such micro-decisions has been a primary driver of progress in the SAT Competition~\cite{satcompetition}, yet each decision is expensive to discover and laborious to validate.   % Recent work has begun to automate this process by coupling Large Language Models (LLMs) with evolutionary search loops, enabling the model to propose, compile, evaluate, and refine heuristic code directly~\cite{sun2024autosat,chen2025dasathco,sun2025automatically,sheng2025solsearch,ding2025selfoptimizing}. % However, these approaches are broadly hindered by two shared and fundamental limitations:  % % We argue that these two capabilities are essential for treating heuristic discovery as a rigorous scientific problem, and that both are missing from prior LLM-based approaches. % %  % The first is \textbf{scale}: meaningful coverage of the design space requires evolving many heuristic functions across multiple solver families in parallel, instead of single-solver, single-function confinement. % Without a throughput sufficient to make systematic study tractable, the search either revisits well-trodden ideas or terminates before non-obvious variants surface. % %  % All existing LLM-based approaches confine their search to a single solver or a narrow slice of the heuristic function space, foregoing the breadth needed to cover the true design landscape.  % This single-solver limitation is in direct tension with a well-established empirical reality: no solver dominates uniformly across all problem categories. % %  % The annual competition confirms this structurally: the winning solver changes every year — Kissat\_MAB-HyWalk in 2022, SBVA\_CaDiCaL in 2023, kissat-sc2024 in 2024 — built on different base lineages and heuristic combinations. % %  % %  % %  % The second is \textbf{controllability}, and the limitation is using the aggregated PAR-2 as the only optimization target, with the absence of controllability over which problem subcategory is targeted. % Practical SAT workloads are often dominated by specific problem structures, and heuristics that excel on aggregate benchmarks may underperform on these subsets. A useful discovery system should be steerable toward such targets rather than producing only one-size-fits-all variants. % %  % However, all existing frameworks optimize a single aggregate PAR-2 score, collapsing SAT, UNSAT, easy, and hard instances into an undifferentiated pool. % %  % This conflation is at odds with how the SAT community evaluates and develops solvers in practice. % The SAT community has long treated these as separate objectives. % SAT Competition evaluations report SAT and UNSAT instance counts independently. %    Recent work has begun to automate this process by coupling Large Language Models (LLMs) with evolutionary search loops, enabling the model to propose, compile, evaluate, and refine heuristic code directly~\cite{sun2024autosat,chen2025dasathco,sun2025automatically,sheng2025solsearch,ding2025selfoptimizing}. However, these approaches are broadly hindered by two shared and fundamental limitations. %  The first limitation is \textbf{scalability}. All existing LLM-based approaches confine their search to a single solver or a narrow slice of the heuristic function space, foregoing the breadth needed to cover the true design landscape. This single-solver confinement is in direct tension with a well-established empirical reality: no solver dominates uniformly across all problem categories~\cite{sheng2025solsearch}. The annual SAT competition confirms this structurally: the winning solver changes every year, built on different base lineages and heuristic combinations~\cite{balyo2022sat,balyo2023proceedings,heule2024proceedings,codel2025sat}. %  The second limitation is \textbf{controllability}. All existing frameworks optimize a single aggregate PAR-2 score, collapsing SAT, UNSAT, easy, and hard test instances into an undifferentiated pool. This conflation is at odds with how the SAT community evaluates and develops solvers in practice: competition evaluations report SAT and UNSAT performance independently~\cite{balyo2022sat,balyo2023proceedings,heule2024proceedings,codel2025sat}. % , and practical scenarios (hardware verification, cryptanalysis, combinatorial planning) are each dominated by specific problem subcategories that aggregate metrics systematically obscure.     % % Two threads of work have sought to reduce this burden. Automated algorithm configuration tools, such as SMAC and irace~\cite{}, tune \emph{existing} heuristic parameters but cannot synthesize new heuristic logic. More recently, large language models (LLMs) have emerged as program synthesizers for algorithmic components~\cite{}, opening the door to automated discovery of new heuristic code rather than mere parameter tuning. Existing LLM-based efforts in this direction, however, remain narrow in scope: they typically target one or two heuristic functions in a single solver, run a small number of generations, and report aggregate improvements without dissecting their source. As a result, they cover only a sliver of the design space and offer no mechanism to steer exploration toward heuristics that excel on particular instance families.     These two bottlenecks lead us to two fundamental questions, i.e., \textit{not how to design yet another LLM-based search algorithm}, but: % These observations motivate two central research questions: \vspace{-0.5em} \begin{center} \fbox{     \parbox{0.9\linewidth}{         % \textit{\textbf{Q1:} Can we discover better SAT heuristics by evolving solvers at \textcolor{red}{large scale}, across many solver families and target functions in parallel?} \\         \textit{\textbf{Q1:} Can we discover better SAT heuristics by evolving solvers at a \textbf{large scale}, in parallel across many solver families and target heuristic functions?} \\         % \textit{\textbf{Q2:} Can we make LLM-driven exploration controllable, so that the discovered heuristics can be specialized to targeted problem subclasses?}         \textit{\textbf{Q2:} Can we make LLM-driven heuristic exploration \textbf{controllable}, such that discovered heuristics can be specialized toward targeted subcategories (SAT, UNSAT, easy, or hard instances) with only a modest trade-off in aggregate PAR-2?}     } } \end{center}  % % KissatEvolve answers these questions through three contributions. \textbf{(i) Controlled mutations} convert retrieved exemplars and structured prompts into precise local moves in heuristic space, replacing the diffuse, overlapping variants typical of unconstrained LLM generation. \textbf{(ii) A scalable asynchronous pipeline} parallelizes solver compilation, execution, and evaluation, enabling the framework to evolve substantially more solvers per unit of wall-clock time than prior systems. \textbf{(iii) A densely annotated memory bank} records both high- and low-performing variants together with their solver analyses, search statistics, and code diffs, converting previous generations into reusable optimization experience that is selectively retrieved to condition future generations. % %  % We instantiate these two principles in a framework, \textbf{KissatEvolve}, that couples large-scale parallel evolution with LLM-based heuristics generation. To realize this design, we confront two technical bottlenecks that have limited prior LLM-driven evolutionary frameworks. They typically lack a structured memory of past trials, causing the search to revisit similar low-value ideas; and absent explicit mutation control, generated variants tend to differ only along vague or overlapping axes, blunting the search signal. As the candidate population grows, orchestration, solver building, and evaluation become further serious system bottlenecks. % % % %  % Together, these components transform LLM-based heuristic discovery from a single-generation search loop into a scalable, retrieval-augmented evolutionary system, and to our knowledge, constitute the first framework to deliver both scalability and controllability for LLM-generated SAT heuristics.   We answer both questions through \textbf{KissatEvolve}, a framework that couples large-scale parallel evolution with retrieval-augmented LLM generation for SAT heuristic discovery. %  % To realize this design, we address two technical bottlenecks that have limited prior LLM-driven evolutionary frameworks: % they typically lack a structured memory of past trials, causing the search to revisit similar low-value ideas; % and without explicit mutation control, generated variants tend to differ only along vague or overlapping axes, blunting the search signal. % As the candidate population grows, solver orchestration, compilation, and evaluation become further system bottlenecks. %  KissatEvolve achieves this through two components:  % \textbf{(i) Controlled mutations} convert retrieved exemplars and structured prompts into precise local moves in heuristic space, replacing the diffuse, overlapping variants of unconstrained LLM generation. \textbf{(i) An automated and scalable pipeline for heuristics generation} that parallelizes solver compilation, execution, and evaluation across multiple solver families simultaneously, enabling the framework to evolve substantially more solvers per unit of wall-clock time than prior systems. \textbf{(ii) Controllability via a densely annotated memory bank} that records both high- and low-performing variants together with solver analyses, search statistics, and code diffs, converting previous generations into reusable experience that is selectively retrieved to steer LLM's generation toward SAT, UNSAT, easy, or hard instance subcategories, allowing researchers to explicitly trade a small amount of aggregate PAR-2 for targeted subcategory gains. %  % is introduced as a first-class objective: % the LLM prompt and fitness signal can be  To our knowledge, this framework constitutes the first framework to deliver both scalability and controllability for LLM-generated SAT heuristics. %  % In the operational mechanism of CDCL SAT solvers, the selection of heuristic strategies directly determines system performance. However, traditional heuristic methods inherently suffer from limited domain adaptability, making it challenging to construct universal solutions. Previous iterations of the AE framework addressed this by leveraging a Mixture of Experts (MoE) approach—integrating models like DeepSeek-R1, Claude 3.7, and ChatGPT-4.5—to drive collaborative learning and genetic evolution. % While the baseline AE framework successfully optimized specific heuristic targets, scaling the evolutionary process required overcoming premature convergence and computational bottlenecks. This work expands the framework into a massive, parallelized pipeline supported by an experience-driven memory structure, enabling the simultaneous evolution of complex solvers including Kissat-Cure, Kissat-Corephase, and Kissat_pred. %  We summarize our contributions below: {\color{red} \begin{enumerate}[leftmargin=*]     \item Using KissatEvolve, we conducted a comprehensive study spanning 25 functions across 6 SAT solvers, which is the first time of large-scale LLM-based heuristics generation.     \item KissatEvolve discovered 5 promising heuristics, improving the PAR-2 scores by XXX\% over state-of-the-art baselines.     \item When controllability is enabled, KissatEvolve achieved 43.76\% PAR-2 improvement on targeted easy, and hard instance subsets. \end{enumerate} }    \section{Background, Motivations, and Challenges}     \subsection{Background: Boolean Satisfiability (SAT) Problems, SAT Solvers, and SAT Heuristics}   % SAT (Satisfiability Problems) is the problem of determining whether there exists an interpretation that satisfies a given Boolean formula. In other words, it asks whether the variables within a given Boolean formula can be assigned values of TRUE or FALSE such that the entire formula evaluates to TRUE.   The Boolean Satisfiability Problem (SAT), the first proven NP-complete problem, asks whether a given propositional formula in conjunctive normal form (CNF, a conjunction of clauses, each a disjunction of literals) can be made true by some assignment of its Boolean variables. A SAT solver takes a CNF instance as input and returns either a satisfying assignment (SAT) or a proof of unsatisfiability (UNSAT). Modern solvers are predominantly based on Conflict-Driven Clause Learning (CDCL), with Kissat~\cite{biere2025cadical} representing the state of the art. The CDCL workflow interleaves several heuristic-governed components, including branching, restart and rephase policies, and clause reduction, each exposing one or more heuristic functions whose tuning critically affects performance. Solver performance is evaluated via the PAR-2 score: the average runtime over a benchmark set, where instances not solved within a timeout $T$ are penalized at $2T$. Instances are further categorized as SAT or UNSAT based on satisfiability, and as easy or hard based on whether contemporary solvers resolve them well within the timeout.       % \textcolor{red}{definition of PAR-2}  % \textcolor{red}{definition of SAT / UNSAT / easy / hard}  % {\color{blue}  % \paragraph{Preliminary Definition on SAT} % % \label{pre:sat} % Let $V = \{x_1, x_2, \dots, x_n\}$ be a set of Boolean variables, a \textit{literal} is either a variable $x$ or its negation $\neg x$. A \textit{clause} is a disjunction of literals. A \textit{conjunctive normal form} (CNF) formula $F = C_1 \land C_2 \land \dots \land C_m$ is a conjunction of clauses. For simplicity we assume all clauses are non-tautological, in other words, no variable $x$ occurs positively $(x \in C)$ and negatively $(\neg x \in C)$ in the same clause.  % A (partial) mapping $\alpha : V \to \{0, 1\}$ is called an \textit{assignment}. If $\alpha$ maps all variables to a boolean value, it is termed a \textit{complete} assignment; otherwise, it is referred to as a \textit{partial} assignment. The value of a variable $x$ under an assignment $\alpha$ is denoted as $\alpha[x]$. An assignment $\alpha$ satisfies a clause if at least one literal evaluates to true under $\alpha$, and satisfies a CNF formula if it satisfies all its clauses. A CNF formula $F$ is satisfiable if there is at least one satisfying assignment. The empty clause $\bot$ is always unsatisfiable and represents a \textit{conflict}. SAT is the problem of deciding whether a given CNF formula is satisfiable.   % \paragraph{CDCL solver} % The CDCL solver is composed of a propagate-and-learn framework and a guessing part. CDCL solvers do propagate-and-learn eagerly in practice and the implementations do not differ much among different solvers. However, the guessing policy and the search-restart policy are important for performance and differ across implementations. Due to the need for manual design and the lack of rigorous mathematical proof, these policies are called heuristics in this paper.   % There is a long history of research on heuristics in SAT solvers, among which branching heuristics play a crucial role and continue to significantly impact the performance of CDCL solvers. For example, Variable State Independent Decaying Sum (VSIDS)~\shortcite{vsids} is a family of branching heuristics that assists in selecting the most promising variable to assign a value in the Make Decision phase. Another important branching heuristic is Learning-Rate Branching (LRB)~\shortcite{lrb}, which frames branching as an optimization problem that picks a variable to maximize a metric called learning rate.  % Restart heuristics is also essential for enhancing the performance of CDCL solvers. It allows the solver to abandon the current search path and backtrack to a specific decision level, but the learned clauses are usually maintained for the following search. Fast restart~\shortcite{ramos2011between} is now widely used. Originally Luby restarts were heavily used because they represent a priori optimal strategy~\shortcite{luby}. However, most implementations switched to Glucose-style restarts~\shortcite{glucose}, which are essentially a requirement for prevailing in the SAT Competition. For an extensive overview of these techniques, see the survey presented in~\shortciteA{biere2015evaluating}.   % Rephase is another technique that introduces heuristics into CDCL solvers. Its primary objective is to reset or adjust the current partial assignment, thereby enabling the solver to explore diverse search space. PrecoSAT and PicoSAT~\shortcite{precosat} utilize a Jeroslow-Wang score~\shortcite{jeroslow1990solving} to adjust the saved phases either on all or only on irreundant clauses in regular intervals, following a Luby sequence. StrangeNight~\shortcite{Strangenight} employs a strategy of flipping values with a certain probability that depends on the depth of the assignment. The motivation is to avoid the heavy-tail phenomenon. What is more, negative results of the SAT solver Riss are reported in~\shortciteA{balint2015overview}, and this is the first time that several rephasing heuristics were compared and used.  % }      % Program synthesis is in general an undecidable search problem~\cite{}. Even with practical assumptions to limit the search, it is often exponential in the number of functions and their size in lines of code~\cite{}. Evolutionary search techniques that combine genetic methods with LLM-based agents, such as AlphaEvolve, have shown great promise in efficiently searching through large search spaces.   % Having said that, there are several key challenges that remain. First, as previously mentioned, the problem requires search through an exponential search space. Second, current open-source versions of AlphaEvolve or other LLM-based evolutionary search methods specifically aimed at searching for highly complex solver heuristics do not take advantage of learning from previous failed iterations. More precisely, when a generated heuristic fails, we can get a lot of rich semantic information from the solver as to why the heuristic failed. Using this information effectively can help prune the search exponentially. However, designing a technique that uses such rich semantic feedback is hard. Third, existing techniques search for single heuristics, limiting the kinds of heuristics one can find. It would be better to search for a set of heuristics that work symbiotically together to make the resultant SAT solver far more efficient.    % These challenges motivate us to look for techniques that can leverage the rich semantic information in clever ways. ....   Automatically synthesizing effective SAT heuristics is an instance of program synthesis over an exponentially large search space: the number of candidate implementations grows with both the number of target functions and their code complexity. Evolutionary search methods that couple genetic operators with LLM-based generation, exemplified by AlphaEvolve~\cite{novikov2025alphaevolve}, have demonstrated strong potential for navigating such spaces efficiently. Yet two fundamental challenges remain underaddressed when applying these methods to SAT heuristic discovery: the search space must be explored \emph{at scale} across multiple solver families and function spaces jointly, and the optimization objective must be \emph{controllable} to target specific problem subcategories. We motivate both requirements below.    \subsection{Improving SAT Solvers Demands Exploration over a Large Heuristics Space.}  % SAT solvers, even just for Kissat, have many variants with different heuristics (i.e., a large space of heuristics). % Evolving just a single solver cannot provide principled/systematic improvements or understanding. % %  % Top groups simultaneously submit multiple solvers from distinct families precisely because cross-family exploration is necessary:  % the ESA group at SC 2024 submits \texttt{CaDiCaL\_ESA}, \texttt{Kissat\_MAB\_ESA}, \texttt{Kissat\_MAB\_Binary}, and \texttt{AMSAT} in parallel \cite{sc2024}. % %  % The cost of staying within a single solver is quantified directly by the ablation in DaSAThco~\cite{chen2025dasathco}: restricting to ``Single Best Selection'' yields PAR-2 = 1711.1, versus 1671.0 for the full search spanning multiple heuristics.  The heuristic design space for SAT solvers is vast: even within a single solver family such as Kissat, dozens of variants exist, each instantiating distinct branching, restart, rephase, and clause-management policies. Exhaustive improvement from a single solver is therefore neither principled nor systematic. This is reflected directly in competition practice: top groups simultaneously submit solvers from multiple distinct families precisely because cross-family exploration is necessary~\cite{heule2024proceedings}. % : the ESA group at SC 2024 submits \texttt{CaDiCaL_ESA}, \texttt{Kissat_MAB_ESA}, \texttt{Kissat_MAB_Binary}, and \texttt{AMSAT} in parallel~\cite{sc2024}. The suboptimal performance of restricting exploration to a single solver is further quantified by the ablation study in DaSAThco~\cite{chen2025dasathco}: a ``Single Best Selection'' policy achieves PAR-2 = 1711.1, compared to 1671.0 for the full portfolio that spans multiple heuristic configurations, a 2.4\% degradation attributable solely to the narrowed search space.   %\begin{wrapfigure}{r}{0.42\textwidth} %    \centering %    \vspace{-0.8em} %    \includegraphics[width=0.4\textwidth]{images/venn_heuristics.pdf} %    \caption{Comprehensively improving SAT solvers requires exploring a broad space of heuristics. In this Venn diagram, three solvers derived from Kissat~\cite{biere2025cadical} all incorporate unique heuristics.} %    \label{fig:venn_heuristics} %    \vspace{-1em} %\end{wrapfigure}  \begin{figure}[h!] \centering \begin{subfigure}[t]{0.32\textwidth}     \centering     \includegraphics[width=\textwidth]{images/venn_heuristics.pdf}     \caption{Three Kissat-derived solvers each contribute unique heuristics. No one solver subsumes the others.}     \label{fig:venn_heuristics} \end{subfigure} \hfill \begin{subfigure}[t]{0.32\textwidth}     \centering     \includegraphics[width=\textwidth]{images/par2_corr_all_ablations_sat_vs_unsat.pdf}     \caption{PAR-2 on SAT vs.\ UNSAT.}     \label{fig:multi_objective_sat} \end{subfigure} \hfill \begin{subfigure}[t]{0.32\textwidth}     \centering     \includegraphics[width=\textwidth]{images/par2_corr_all_ablations_easy_vs_hard.pdf}     \caption{PAR-2 on easy vs.\ hard.}     \label{fig:multi_objective_hard} \end{subfigure} \caption{\textbf{Two motivations for KissatEvolve.} (a) Improving SAT solvers requires exploring heuristics across solver families. (b, c) Weak SAT/UNSAT and easy/hard correlations show that heuristic optimization is multi-objective.} \label{fig:motivations} \end{figure}  We use a Venn diagram to further justify this gap. As illustrated in Figure~\ref{fig:venn_heuristics}, three representative solvers submitted to recent SAT competitions share a common implementation core, yet each still introduces heuristics absent from the others, and no single solver subsumes the heuristics of the rest. Comprehensive improvement therefore cannot be achieved by refining one solver in isolation: it requires exploring a heuristic space that spans the union of these disjoint regions, motivating frameworks that search across solvers and heuristic families jointly. %  Despite all of this evidence, no existing LLM-based framework has been designed to jointly evolve heuristics across multiple solver families and function spaces at scale, leaving the vast majority of the design space systematically unexplored.      % \textcolor{red}{Todo @Shimin: show a Venn diagram of shared/unique functions in Kissat variants?}  \subsection{Heuristics Optimization is Multi-Objective} \label{sec:motivation_multiobjective}  % \textcolor{red}{Todo @Meru/Shimin: Show a diagram (bar plots: 1) y-part2, x-SAT/UNSAT; 2) y-part2, x-easy/hard) of imbalanced par-2 over SAT/UNSAT, easy/hard instances.}  % \begin{figure}[h!] % 	\centering %     % \captionsetup{font=small} % % \vspace{-0.5 em} % 	\caption{\textbf{Weak correlations} (PAR-2 of SAT vs. UNSAT) indicate the nature of being multi-objective when improving SAT heuristics.} % \label{fig:multi_objective2} % % \vspace{-1.5 em} % \end{figure}   % In practice, % SAT community evaluates solvers' SAT and UNSAT performance independently~\cite{balyo2022sat,balyo2023proceedings,heule2024proceedings,codel2025sat}, and historically solvers' rankings on SAT and UNSAT benchmarks are typically not aligned. % To further justify this multi-objective nature of SAT heuristics optimization, % we plot across different solvers and heuristics (generated by our LLM in experiments, see Section~\ref{sec:experiment}). % %  % As shown in Figure~\ref{fig:multi_objective}, PAR-2 of solvers on sub-categoriy problems (SAT vs. UNSAT; easy vs. hard thresholded by PAR-2 = 100) are not aligned, indicating that optimizing the performance on each category requires controllable exploration over heuristics. % %  % Yet no existing LLM-based discovery system exposes a mechanism to steer generation toward any of these targets: practitioners who require category-specific improvements have no recourse but to return to manual solver engineering, precisely the bottleneck this line of research was meant to address. % %    In standard SAT competition evaluation, solver performance on satisfiable and unsatisfiable instances is assessed independently~\cite{balyo2022sat,balyo2023proceedings,heule2024proceedings,codel2025sat}. Rankings on SAT versus UNSAT benchmarks are historically misaligned: a solver that excels on one category frequently underperforms on the other.  To make this multi-objective structure concrete, Figure~\ref{fig:multi_objective} plots pairwise PAR-2 scores across a diverse set of solver configurations generated by our LLM-based framework (see Section~\ref{sec:experiment}), partitioning instances into four sub-categories: SAT vs.\ UNSAT, and easy vs.\ hard (thresholded at PAR-2 = 100). The relatively low correlation between PAR-2 on each pair of sub-categories demonstrates that heuristic choices which improve performance in one regime may fail to transfer to others, optimizing for SAT does not yield commensurate gains on UNSAT, and vice versa. This Pareto-style tension implies that SAT heuristics optimization is inherently multi-objective, and that controllable, category-aware exploration of the heuristic space is necessary to make targeted progress. %  Yet no existing LLM-based heuristic discovery framework exposes a steering mechanism for any of these objectives: practitioners who require improvements on a specific sub-category are forced to resort to manual solver engineering, precisely the bottleneck that automated heuristic optimization is intended to eliminate.       \section{KissatEvolve: Controllable and Scalable Synthesis of SAT Heuristics}   \textcolor{red}{need to \\ref Table 1 in main text.}  \textcolor{red}{Inconsistency: 1. parent vs leader: use parent, 2. naming format of solver and functions: font, hyphen or underscore.}  \vspace{-0.5em} \begin{table}[H] \centering \caption{The scale of our KissatEvolve outperforms previous LLM-based works on the number of studied solvers, heuristics, and controllability.} \label{tab:comparison} \resizebox{0.63\linewidth}{!}{% \begin{tabular}{lccc} \toprule Method & \#Solvers & \#Heuristics & Controllability \\ \midrule AutoSAT~\cite{sun2024autosat} & 1 (EasySAT) & 9 & No \\ DaSAThco~\cite{chen2025dasathco} & 1 (EasySAT) & 3 & No \\ AutoModSAT~\cite{sun2025automatically} & 1 (ModSAT) & 7 & No \\ % SolSearch & ? & ? & X \\ AE Framework~\cite{ding2025selfoptimizing} & 1 (Kissat) & 3 & No \\ % KissatEvolve (Ours) & \textbf{6} (Kissat variants) & \textbf{25} & Yes \\ KissatEvolve (Ours) & \textbf{6} (Kissat variants) & \textbf{16} & Yes \\ \bottomrule \end{tabular} } \end{table}     Our KissatEvolve operates at the heuristics (function) level. %  For each target solver and target heuristic, KissatEvolve executes the following steps: 1) prompts an LLM to propose (generate) a natural-language algorithmic modification, 2) translates that proposal into code, 3) injects the generated code into the corresponding solver function, 4) builds the solver, 5) validates the result, and 6) evaluates performance. %  % The evolution stages are: parent generation, mutation, and crossover. Core stages for the generation (step 1) include parent generation and mutation (Section~\ref{sec:diversified_mutation}), connected through a high-throughput execution pipeline (Section~\ref{sec:high_performance}) and a persistent memory bank (Section~\ref{sec:memory_bank}).  % During evaluation, the system extracts generated code from model output, replaces the exact target function body using a registry of function locations, and builds solver copies automatically. %  % Quick evaluation jobs are then submitted in batch mode, % after which the best-performing team members are promoted to leaders. %  A validation stage also checks solving logs and proof status so that invalid solver outputs can be detected and excluded before promotion or reuse  (see Appendix~\ref{app:step_decomposition} for detail). %  This design allows natural-language ideas, code implementations, and evaluation signals to be linked end-to-end in a reproducible optimization loop.   \begin{figure}[h!] 	\centering 	\includegraphics[width=0.6\textwidth]{images/overview_mutation.pdf}     % \captionsetup{font=small} % \vspace{-0.5 em} 	\caption{Overview of our framework:     1) Parent generation with diversified mutation (Sec.~\ref{sec:diversified_mutation}),     2) High-Performance Parallel Evolution (Sec.~\ref{sec:high_performance}),     3) Densely Annotated Memory Bank (Sec.~\ref{sec:memory_bank}).     } \label{fig:overview} % \vspace{-1.5 em} \end{figure}    % \textcolor{red}{Need a (sub)section in appendix to briefly discuss and give some numbered evidence why we believe the cross-over stage may not be very help.}   \subsection{Scalable and Automated Evolution of SAT Heuristics with Diversified Mutation}     \subsubsection{Parent Generation with Diversified Mutation} \label{sec:diversified_mutation}   % \textcolor{red}{Todo @Meru @Mike: details on leader generation and mutation.} \\ % \textcolor{red}{Mike: I've added how mutation pool is used in the mutation process}  % \textcolor{red}{Meru: I've added a bunch of details and did some touch up.}  \paragraph{Aim: Generating diverse variants (parents and mutants) of SAT heuristics in natural language.} In this step, we demonstrate our approach to generate heuristics in natural language and achieve high diversity in output to tackle the largest possible search space.     \paragraph{Why is promoting diverse heuristic generation important?}  % !!! To promote diversity across parents within a single iteration, we linearly vary the sampling temperature from 0.5 to 1.0 over the batch to promote different levels of exploration. We further protect against heuristic collapse with parent monitoring using pairwise cosine similarity (Qwen-3 embedding model) and an LLM-Judge that flags strategies which are mechanistically equivalent.   \paragraph{Overview of Steps.} \ul{Input: Parent algorithm generator prompt}; \ul{Step 1: Use mutant generator prompt to generate mutants of each parent algorithm}; \ul{Step 2: Evaluate all parent and mutant algorithms to select the best from each team}; \ul{Output: Re-select highest performing per-team as parents and save all results to memory bank.}    The following is iterated three times:   \textbf{Step 1: Parent Generation.} \enspace Each iteration begins by generating a population of \textit{parent algorithms}: high-level heuristic strategies expressed in \textit{natural language}. We query an LLM with the target heuristic function and its current implementation in Kissat and instruct it to propose a mechanistically distinct strategy for that function. %  The output specification of each parent is a JSON containing both an algorithm and a justification for the improvements of that algorithm  (see Appendix~\ref{app:step_decomposition} for detailed implementation and Appendix~\ref{app:leader_prompt} for LLM prompt).      \textbf{Step 2: Mutant Generation.} \enspace Our contribution is diversified mutations. %  Instead of asking the LLM to generate loosely related variants of a parent. The key idea is to require the parent description to be written as an explicit sequence of algorithmic steps. The system parses these steps, determines how many mutation opportunities exist, assigns generated variants to specific steps, and instructs the LLM to modify only the assigned step for each mutant. We refer to Appendix~\ref{app:mutant_prompt} for the mutant prompt. In this way, mutant generation becomes a structured local search around the parent algorithm rather than an unconstrained rewrite.   This design has two advantages. First, it improves diversity in a meaningful way: mutants differ along interpretable axes rather than through noisy global edits, while the retrieved exemplars steer each step's modification toward locally-validated trajectories. Second, because each mutant is tied to a particular mutated step, the same key supports both storage and retrieval, turning diversified mutation into a self-improving local search whose distribution over edits is shaped by accumulated step-level evidence rather than by a single generic prompt.    % To effectively explore local heuristic spaces without losing the core logic of promising algorithms, we introduce a controlled mutation mechanism. Rather than allowing the LLM to rewrite entire algorithms indiscriminately, the framework enforces step-wise isolation:  % In the initial team leader prompt, the LLM is instructed to explicitly add numbered "Steps" to its algorithm description.  % The system parses these steps to determine the total number of algorithmic components.  % Based on the number of requested variants, the system assigns specific variants to target different individual steps.  % During variant generation, the LLM is strictly instructed to only mutate the specific step it was assigned, leaving the rest of the algorithm intact.   % Step II - Mutant Generation: For each parent solver, create m variant strategies as offsprings through controlled mutations. % To avoid generated mutations to have large overlaps to each other, we first prompt LLMs to divide the solver function into semantically meaningful stages, and each control each mutation to only rewrite the specific stage, without modifying other stages. % Both parents and mutants will be evaluated with PAR-2 score.    \textbf{Step 3: Parent Re-selection.} \enspace    Between iterations of mutations, mutants with the lowest PAR-2 scores are promoted to parents. The rest of the algorithms are stored into our densely annotated memory bank (see Section~\ref{sec:memory_bank}).      We also designed a cross-over phase of the pipeline, but empirically observed that mutation drives most of the performance gain per iteration, and thus elected to focus only on successive mutation iterations for most of our search time. We refer to Appendix~\ref{app:cross_over_detail} for detailed explanation.  % run it for three rounds before each crossover stage.    % \subsection{Cross-Over} % \label{sec:cross_over}  % \textcolor{red}{Todo @Jialin}  % \textcolor{red}{Jialin: Reminder: need to be consistent with ``leaders vs parents'' for Methods}  % \textcolor{red}{Jialin: 1. Update Figure 3 with step names; 2. Many LLM prompts go into appendix; 3. Score threshold in Step II needs justification in appendix.}  % Taking parents from the mutation stage as inputs, our framework incorporates an independent cross-over process to further enhance the solver by combining complementary strengths across algorithms from various mutation groups.  % \textbf{Strategy Analysis.} \enspace For each parent, the large language model (LLM) receives the full algorithm description, its C implementation, and its PAR-2 score. The LLM is then asked to decompose each parent's behavior into four categories: strengths, weaknesses, key mechanisms, and improvement suggestions. This decomposition forces the model to reason about why an algorithm performs the way it does, and provides principled grounding for the subsequent combination step in identifying which parent strengths are complementary and which weaknesses each can compensate for.  % \textbf{Hybrid Evolution.} \enspace Instead of enumerating all combinations of candidate pairs, KissatEvolve prompts the LLM for pair selection. The LLM receives all leaders with their causal analysis and relevant past combination outcomes retrieved from the memory bank in Sec.~\ref{sec:memory_bank}. It then proposes the top-$k$ most promising pairings, each accompanied by a complementarity score and a verbal justification based on the causal reports. Only proposals above a score threshold (6 out of 10) proceed to crossover. This design has two advantages over exhaustive enumeration: it scales more efficiently with population size, and it selectively chooses the pairs with complementary strengths.  % \textbf{Evaluation and Pruning.} \enspace For each accepted pair, the LLM generates a new offspring algorithm based on the analysis from Steps I and II. The offspring description is required to follow a step-by-step structured format, so that the memory bank (Sec.~\ref{sec:memory_bank}) can later attribute the success or failure of a combination to specific algorithmic decisions. The offspring description is then translated into C code by a separate LLM call using both parent implementations as reference. After evaluation, offspring that fail to improve over the best parent PAR-2 are discarded, ensuring the population only retains beneficial combinations.  % Step I - Strategy Analysis: Employ large language models % (LLMs) to conduct causal analysis of parent strategies, identifying precise mechanisms underlying their performance improvements or % degradations.  % TODO: Cross-Over Memory  % Step II - Hybrid Evolution: Based on the detailed examination of how the model influences outcomes, we employ the LLM to identify the most % promising strategy combinations for experimentation, followed by evaluation and strategic analysis.  % Step III - Evaluation and Pruning: Retain the top-performing strategy as new parents.  \subsubsection{Code Generation} \label{sec:code_generation}  \paragraph{Aim: Implement SAT heuristics from natural language into C.} In this step, our pipeline implements the natural-language algorithms from Section~\ref{sec:diversified_mutation}  in C to facilitate evaluation of novel, promising solvers.     \paragraph{Why is automated and scalable implementation across solvers and heuristics challenging?}  % !!! \enspace We expect the coder model (primarily Gemini 3.1 Pro Preview) to have substantial prior exposure to standard heuristic functions in Kissat, given the solver's prominence and documentation. However, KissatEvolve is designed to \textbf{zero-shot transfer to a wider range of solvers and on functions} not present in the original Kissat codebase, which exhibit largely heterogeneous designs, implementations, and macros, incurring significant challenges in seamless adoption of KissatEvolve to different solvers and heuristics.  \textbf{Code Generation.} \enspace  Each natural-language algorithm produced in the parent and mutation stages is passed to a separate coder LLM that implements it in C. After the code is ready, our pipeline copies the base solver directory, locates the point in the code and injects the generated algorithm, and then builds and configures the solver so that the binary is ready for evaluation. We refer to Appendix~\ref{app:code_gen_prompt} for the code generation prompt.   \textbf{Signal Reference.} \enspace To address the codebase awareness concern, we supply the coder with a \textit{signal reference}, which is a markdown sheet enumerating the data structures, types, and solver-internal APIs available to the target function. This reference raised the compilation rate across a team (consisting of a parent algorithm and its mutants) from ~35\% to 90\%. More implementation detail is provided in the Appendix~\ref{app:signal_reference}. This reference was automatically inserted into the coder prompt upon selection of a new target function (see Section ~\ref{sec:target_solvers}) by the pipeline.       \subsubsection{Generating Large-Scale SAT Solver with High-Performance Parallel Evolution} \label{sec:high_performance}  \paragraph{Aim: Highly parallelized and large-scale discovery of SAT heuristics.} In this step, we implement a pipelined architecture for parallel evolution which results in an average 69.41\% decrease in total wall-clock time.     %\paragraph{Why is parallel computing blocking automated and scalable heuristic discovery?}  % !!!     \paragraph{Overview of Steps.} \ul{Input: All algorithm and coder prompts}; \ul{Step 1: Leader and mutant algorithm generation}; \ul{Step 2: Natural-language algorithms deposited into code queue}; \ul{Step 3: Code generation workers pull from the code queue}; \ul{Step 4: Implemented algorithms deposited into build queue}; \ul{Step 5: Build workers pull from build queue and submit evaluations}; \ul{Output: Best performing algorithms re-selected as parents.}   % \paragraph{High-Performance Parallel Evolution}  % \textcolor{red}{Todo @Meru @Jialin: add more details.}   Our second contribution is a high-performance evolution engine that scales the number of solver candidates explored concurrently. LLM-based solver evolution is limited not only by search ideas but also by API throughput, compiler time, job submission overhead, and metadata synchronization; we address these bottlenecks through staged queues, semaphores, and dedicated worker pools, preserving the parent-mutant relationships needed for later promotion and experience updates.   \textbf{Pipeline Design.} \enspace The pipeline is organized as three asynchronous stages connected by two queues: algorithm generation, code generation, and build-and-submit, each producing the input for the next on a per-algorithm basis. Algorithm generation dumps natural-language parents and mutants into the code queue as they are produced; code-generation workers continuously poll this queue, send API requests to the coder model, and enqueue the resulting C implementations onto the build queue. Build-and-submit workers then allocate disk space, inject the generated code into a fresh copy of the base solver, and dispatch an evaluation job array to a compute node without blocking; PAR-2 results return asynchronously via SLURM and drive parent promotion in subsequent iterations. The orchestrator itself runs on the login node and manages concurrency: API calls go through a thread pool gated by the model provider's rate limit, while build-and-submit and database access use separate worker pools so the DynamoDB instance is not overwhelmed at stage boundaries.  % pipeline diagram \textcolor{red}{FIGUREX}? % shows the design with 50 concurrent parent-generation coroutines, 10 variants per parent, 5 asynchronous code-generation workers, and 3 asynchronous build-submit workers.        % To scale the optimization of SAT solvers, the framework utilizes an asynchronous, stage-gated orchestration pipeline that decouples generation, coding, and compilation.  % Stage 0 (Leader Generation): Utilizes up to 50 concurrent coroutines gated by an API semaphore to rapidly generate diverse heuristic leaders.  % Stage 1 & 2 (Variant & Code Generation): Processes 10 variants per leader using concurrent async workers, decoupling natural language algorithm generation from the actual C code synthesis via a code_queue.  % Stage 3 (Build & Submit): Employs 3 async workers operating on a build_queue to compile the solvers (./configure && make) locally, followed by dispatching validation runs via SLURM job arrays (sbatch).        \subsection{Controllable Heuristics Generation via Densely Annotated Memory Bank} \label{sec:memory_bank}  % \textcolor{red}{Note that our memory bank enables implicit communication among parents/mutations via experience sharing.}  % \textcolor{red}{Todo @Mike: add more details, plus a figure to show the data structure of our memory bank.} \\ % \textcolor{red}{Mike: Done, added more details about mutation pool. Since the combination process is not used, i didn't include anything about the combination pool}  \begin{figure}[h!] 	\centering 	\includegraphics[width=1\textwidth]{images/memory_bank.pdf}     % \captionsetup{font=small} % \vspace{-0.5 em} 	\caption{To enable controllable heuristics generation, we introduce the memory bank for the mutation step. Queries retrieve top-$K$ GOOD and BAD exemplars via FAISS indexing to ground the Mutation Prompt. The memory bank is iteratively updated based on solver iteration results, retaining records ranked by relative PAR-2 change.     } \label{fig:memory_bank} % \vspace{-1.5 em} \end{figure}  % \textbf{Densely Annotated Memory Bank.} % \enspace % Our third contribution is a densely annotated memory bank, which we call the mutation experience pool,  \paragraph{Aim: Reusable and controllable SAT heuristics discovery.} In this step, we leverage a memory bank design so that the algorithm generator model learns from previously generated heuristics.     %\paragraph{Why do we need to make LLM-based generation of SAT heuristics sustainable and controllable?}  % !!!     \paragraph{Overview of Steps.} \ul{Input: Parent algorithm $\ell$ and target mutation step $s$}; \ul{Step 1: Embed the query ``Parent: $\ell$ / Step: $s$'' and retrieve top-$K$ records from the good and bad FAISS partitions}; \ul{Step 2: Inject examples of "what worked" and "what to avoid" into the mutation prompt}; \ul{Step 3: After mutant evaluation, rank new (parent, mutant) pairs by relative PAR-2 change and write top-$K$ improvers and top-$K$ degraders back to their respective partitions}; \ul{Output: Mutation prompt grounded in useful examples.}       Large-scale generation of SAT heuristics naturally produces many records of past attempts, each carrying reusable information about what kinds of edits work, what kinds fail, and under what context. Rather than discarding this signal or storing only final PAR-2 scores, we adopt a memory bank~\cite{zhong2024memorybank} that stores rich comparative records paired with LLM-generated analyses and exposes them through semantic retrieval, turning the diversified mutation step into experience-driven local search across iterations. Crucially, our scope goes beyond improving overall PAR-2: as discussed in Section~\ref{sec:motivation_multiobjective}, practitioners need to balance overall PAR-2 against sub-category performance, which requires storing fine-grained per-category results and steering generation via filtered retrievals.   \paragraph{Design.} We illustrate our memory bank design in Figure~\ref{fig:memory_bank}. The pool is organized into two \textcolor{red}{FAISS-indexed} partitions, a \emph{good} partition for mutants that lowered PAR-2 below their parents and a \emph{bad} partition for mutants that regressed. Each record stores the parent's step-decomposed description, the specific step that was mutated, the resulting mutant description, parent and mutant PAR-2 scores broken down by sub-category (overall, SAT, UNSAT, hard, easy), and an explanatory analysis (see Appendix~\ref{app:analysis_generation}). Records are embedded with Qwen3-Embedding-0.6B from the labeled key ``Parent Algorithm Description: $\langle\ell\rangle$, Mutation Step: $\langle s\rangle$'', so the same step-aware representation used by mutant generation is reused as both storage and retrieval handle. After every iteration, only the top-$K$ improvers and top-$K$ degraders, ranked by relative PAR-2 change against the parent, are persisted, keeping the bank focused and bounded.   This dense annotation is what distinguishes the bank from a simple archive. At the next mutation step, the system retrieves semantically similar past successes and failures keyed by the current (parent, step) and injects them into the mutation prompt as ``what worked'' and ``what to avoid'' exemplars, each accompanied by its own short analysis. Retrieval can optionally be reranked by a chosen PAR-2 sub-category (SAT, UNSAT, hard, easy) to bias exploration toward a target objective. Search is thus driven not by raw scores alone but by accumulated, code-grounded reasoning about why specific local edits succeeded or failed, making the framework history-aware in a way that compounds across generations.   \textcolor{red}{code-diff?}  % To prevent the evolutionary process from revisiting failed paths and to ensure continuous learning across iterations, we implemented a densely annotated memory bank. When generating diverse team leaders, providing only "good" examples can cause premature convergence, so the system is explicitly fed information on what constitutes a "bad" strategy. The memory bank is divided into three distinct pools:  % Algorithm Experience Pool: Retains the worst-performing algorithms. It stores the bad algorithm description alongside an LLM-generated analysis detailing exactly why it underperformed compared to the baseline.  % Mutation Experience Pool: Tracks the parent-child relationship between team leaders and their mutated members. It stores the code difference, the specific step that was mutated, the resulting PAR-2 scores, and an analysis of why the mutation improved or degraded performance. This pool is queried using a (leader, step) tuple.  % Combination Experience Pool: Captures the outcomes of genetic crossover events. It stores descriptions of both parent algorithms, the resulting child algorithm, and a causal analysis of why the child outperformed or underperformed its parents.       \section{Experiments} \label{sec:experiment}  \subsection{Settings}   Our primary evaluation set are the 400 CNF instances provided by the SAT competition 2025~\cite{satcomp2025}. In order to facilitate faster evaluation, we selected a subset of 50 CNF instances and lowered the PAR-2 timing to a quick evaluation standard measured with a 600s timeout per instance and a 1200s penalty on timeout. More detail is provided in the Appendix~\ref{app:benchmark_subset}.  We primarily used both gemini-3.1-pro-preview and gemini-3-flash-preview as the algorithm generation as well as the coder model due to its ability to achieve high compilation rates at a low-cost.  Evaluations were split between two clusters. Most ran on x86\_64 dual-socket AMD EPYC (Zen 5/Turin) nodes (192 cores/768 GB RAM per node, no SMT). Another cluster contains 2 AMD EPYC 7763 (Milan) CPUs with 64 cores each and 512 GB RAM in total. Each solver built on a single core under an 8 GB memory limit, scheduled via SLURM job arrays.    \subsection{Target Solver Families and Target Functions} \label{sec:target_solvers}  % \textcolor{red}{Todo @Shimin: please provide as many as possible justifications on why choosing these solvers / functions. We need a structured way to organize these choices.}  Kissat\cite{biere2025cadical} is one of the state-of-the-art SAT solvers, and many solvers in the SAT competition\cite{satcomp2025} are variants of Kissat. We targeted the best-performing Kissat variants from the SAT competition 2025~\cite{satcomp2025}, aiming to build a competition-winning solver.  From the top Kissat variants, we selected six solvers based on the following criteria: lower PAR-2 score preferred, and base solver if many variants exist.   For AE-Kissat-MAB\cite{ding2025selfoptimizing}, the winner from SAT competition 2025~\cite{satcomp2025}, we generated functions for the restart policy and branching heuristic, two of the most influential heuristics in a CDCL SAT solver. For other solvers, we performed generation on a multitude of heuristics to test KissatEvolve's capability on improving less commonly studied functions. A more thorough explanation of heuristic choice is available in the Appendix.      % Experimental Targets: Large-Scale SAT Solvers % The enhanced framework was deployed to evaluate and optimize a diverse set of functions across multiple large-scale SAT solver variants. The following solvers and their respective core functions were targeted for evolution:   % \textcolor{red}{Todo: we may use LLM to identify these heuristics, beyond the scope of this work. Even we manually picked these heuristics, zero-shot transfer KissatEvolve to different heuristics is still non-trivial.}   % Kissat\_vsa\cite{li2025kissatvsa}: Optimized functions include init\_activityref, rank\_vivification\_candidates\_by\_activity, and sort\_vivification\_candidates\_by\_activity.   % Kissat\_pred\cite{luo2025kissatpred}: Evolution focused on predictive and configuration parameters: kissat\_set\_configuration, predict, extract\_features, and extract\_sequence\_features.   % Kissat-Corephase\cite{chen2025kissatcorephase} \& Kissat\_Corephase\_Coreward: Targeted complex phase-saving and walking behaviors, including reset\_phases\_mab, walking\_phase, rephase\_conflict, largest\_score\_unassigned\_variable, kissat\_analyze, rephase\_walking, reset\_phases, and kissat\_rephase.   % Kissat-Cure\cite{cai2025kissatcure}: Evolved functions related to dynamic phase decisions and restarts, specifically kissat\_dynamicsat, kissat\_decide\_phase, kissat\_reboot\_direct, and kissat\_coldrestarting.   % AE\_kissat2025\_MAB: Optimized restarting and scoring mechanics via kissat\_restarting, restart\_mab, kissat\_bump\_score\_increment, and kissat\_rescale\_scores.      \subsection{Large-Scale Synthesis of SAT Heuristics}   % We first preliminarily study all heuristics with selected \textcolor{red}{40} problems. % As shown in Table~\ref{tab:solver_comparison_quick}, we successfully identify candidates that are promising to outperform each of their baselines.  We begin with a preliminary screening of all generated heuristics on a selected subset of \textcolor{red}{40} problems to efficiently allocate evaluation budget. As shown in Table~\ref{tab:solver_comparison_quick}, this screening successfully identifies candidates that surpass their respective baselines: notably, \texttt{kissat\_dynamicsat} (469.44 vs.\ baseline 531.05), \texttt{kissat\_bump\_score\_increment} (475.58 vs.\ baseline 528.76), and \texttt{rephase\_conflict} (603.94 vs.\ baseline 740.72).    % We thus allocate more \textcolor{red}{(how many? roughly)} budget to these promising heuristics and generate more parents, evaluating with full testing set. % We identified some heuristics that can significantly outperform the baseline, especially ``kissat\_dynamicsat'' of solver ``Kissat-Cure'' and ``kissat\_bump\_score\_increment'' of ``AE\_kissat2025\_MAB''  We then allocate additional computational budget to these promising heuristics and evaluate them on the full benchmark set. As shown in Table~\ref{tab:solver_comparison}, the preliminary rankings prove predictive: both \texttt{kissat\_dynamicsat} (Kissat-Cure) and \texttt{kissat\_bump\_score\_increment} (AE\_kissat2025\_MAB) substantially outperform their baselines, with PAR-2 reductions of 35.4\% (2345.57 vs.\ 3630.45) and 1.5\% (1859.38 vs.\ 1887.43), respectively. The improvement of \texttt{kissat\_dynamicsat} is particularly striking, suggesting that dynamic activity-based scheduling offers a fundamentally more effective strategy for this solver. \texttt{rephase\_conflict} also improves over its baseline, though more modestly (1.5\%), indicating that phase-reset heuristics may offer narrower gains for Kissat\_Corephase\_Coreward. See Appendix~\ref{app:examples} for the full natural-language descriptions  and C implementations of representative generated heuristics.   \textcolor{red}{Todo: much more detailed analysis on what these two heuristics are doing and why improved.}    \begin{table}[h] \centering \captionsetup{font=small} \caption{Final Solver comparison results.} % \vspace{-1.0em} \resizebox{.9\textwidth}{!}{ \begin{tabular}{|l|l|c|c|} \hline \textbf{Solver} & \textbf{Heuristics} & \textbf{Best PAR-2 $\downarrow$} & \textbf{Base PAR-2 $\downarrow$} \\ \hline \multirow{1}{*}{Kissat-Cure}   & kissat\_dynamicsat  & 2345.57 & 3630.45 \\   % & kissat\_decide\_phase & 2315.53 & 2206.33 \\ \hline Kissat\_Corephase\_Coreward & rephase\_conflict & 2459.43 & 2497.15 \\ \hline \multirow{2}{*}{AE\_kissat2025\_MAB}   & restart\_mab                  & 1922    & \multirow{2}{*}{1887.43} \\   & kissat\_bump\_score\_increment & 1859.38 & \\ % \hline % Kissat-Cure & kissat\_decide\_phase & 2198.71 & 2165.94 \\ \hline \end{tabular} } \label{tab:solver_comparison} \end{table}     \subsection{Memory Bank Improves PAR-2 Scores} To evaluate the effectiveness of the memory bank, we ablate on two heuristic functions from the Kissat-Cure solver: \texttt{kissat\_dynamicsat} and \texttt{kissat\_decide\_phase}. For each configuration, we run one iteration of diversified mutation with 3, 5, 7, 9, and 11 parents, each generating 3 mutants, and repeat each setting 3 times to reduce variance. The memory bank retrieves at most 3 good and 3 bad exemplars per query. As shown in Figure~\ref{fig:ablation1}, enabling the memory bank improves PAR-2 ($\downarrow$) in the majority of configurations, with peak improvements of 80.05 and 118.53 for \texttt{kissat\_dynamicsat} and \texttt{kissat\_decide\_phase} respectively, demonstrating that conditioning mutation on retrieved experience consistently steers generation toward better-performing heuristics.   \begin{figure}[h]     \centering     \begin{subfigure}{0.49\textwidth}         \includegraphics[width=\linewidth]{images/ablation1_kissat_dynamicsat.pdf}         \caption{kissat\_dynamicsat Heuristics in Kissat-Cure Solver}     \end{subfigure}     \hfill     \begin{subfigure}{0.49\textwidth}         \includegraphics[width=\linewidth]{images/ablation1_kissat_decide_phase.pdf}         \caption{kissat\_decide\_phase Heuristics in Kissat-Cure Solver}     \end{subfigure}     \caption{Effect of memory bank and parent count on PAR-2. \textbf{Memory on}: memory bank is used to retrieve exemplars during mutation; \textbf{Memory off}: no retrieval is used. PAR-2 is evaluated on the 50-instance benchmark subset (Appendix~\ref{app:benchmark_subset}). Error bars denote standard deviation over 3 runs.}     \label{fig:ablation1} \end{figure} % (w/ vs. w/o memory bank) better performance / sampling efficiency  % X: \#leaders (4~5 \# of leaders); Y: Par2; two curves    \subsection{Controllability via Memory Bank}  % Table: two rows (baseline “original memory bank” vs. controlled); four columns of PAR-2: sat/unsat/easy/hard  To evaluate controllability, we compare two retrieval strategies: \textit{controlled retrieval}, which reranks the top-10 retrieved exemplars by their PAR-2 score on the target subcategory (Easy or Hard) before selecting the final 3 good and 3 bad examples, against \textit{uncontrolled retrieval}, which selects exemplars by semantic similarity alone without subcategory re-ranking. Both settings use the same 3, 5, 7, 9, and 11 parent configurations with 3 mutants per parent. As shown in Table~\ref{tab:ablation2}, controlled retrieval improves PAR-2 in the majority of Easy and Hard configurations for both heuristics, confirming that subcategory-aware reranking effectively steers LLM generation toward instance-specific objectives without requiring changes to the generation pipeline itself. % Fixed \#leader \begin{table}[h] \centering \captionsetup{font=small} \caption{PAR-2 ($\downarrow$) of controlled vs.\ uncontrolled memory bank retrieval on Easy and Hard instances. \textbf{w/}: controlled retrieval, where exemplars are re-ranked by subcategory-specific PAR-2 before selection; \textbf{w/o}: uncontrolled retrieval, where exemplars are selected by semantic similarity alone. PAR-2 is evaluated on the 50-instance benchmark subset; Easy/Hard subcategory definitions and subset construction are detailed in Appendix~\ref{app:benchmark_subset}.} \resizebox{\textwidth}{!}{ \begin{tabular}{|l|l|c|c|c|c|c|c|c|c|c|c|} \hline \multirow{2}{*}{\textbf{Heuristics}} & \multirow{2}{*}{\textbf{Split}}   & \multicolumn{2}{c|}{\textbf{Parents=3}}   & \multicolumn{2}{c|}{\textbf{Parents=5}}   & \multicolumn{2}{c|}{\textbf{Parents=7}}   & \multicolumn{2}{c|}{\textbf{Parents=9}}   & \multicolumn{2}{c|}{\textbf{Parents=11}} \\ \cline{3-12} & & \textbf{w/} & \textbf{w/o} & \textbf{w/} & \textbf{w/o} & \textbf{w/} & \textbf{w/o} & \textbf{w/} & \textbf{w/o} & \textbf{w/} & \textbf{w/o} \\ \hline \multirow{2}{*}{\texttt{kissat\_dynamicsat}}   & Easy & \textbf{280.15} & 461.55 & \textbf{256.16} & 455.50 & \textbf{559.60} & 562.42 & \textbf{244.42} & 417.69 & 443.40 & \textbf{439.82} \\   & Hard & \textbf{1108.61} & 1138.14 & \textbf{1157.07} & 1164.52 & \textbf{1142.53} & 1182.69 & \textbf{1129.32} & 1157.00 & \textbf{1137.16} & 1161.24 \\ \hline \multirow{2}{*}{\texttt{kissat\_decide\_phase}}   & Easy & \textbf{126.45} & 210.77 & \textbf{164.88} & 205.44 & \textbf{186.93} & 263.99 & \textbf{154.53} & 217.47 & 217.24 & \textbf{177.73} \\   & Hard & \textbf{954.73} & 1022.19 & \textbf{998.78} & 1014.38 & \textbf{1001.76} & 1075.69 & \textbf{949.52} & 1053.62 & \textbf{961.20} & 1018.26 \\ \hline \end{tabular} } \label{tab:ablation2} \end{table}  % \subsection{``Trustworthiness'' of LLM-generated heuristics}  % validation of solvers  % % is a heuristic empirically (on benchmark) efficient?  % code diff  % four metrics, LLM-prompted insights  % % theorem of “fake/surrogate mode” of solvers  % % if magic numbers changed: completely fail the solver   \section{Related Works}  code synthesis  heuristics discover  ML for Solvers    \section{Limitation}    \section{Conclusion}    % \section{Submission of papers to NeurIPS 2026}   % Please read the instructions below carefully and follow them faithfully.   % \subsection{Style}   % Papers to be submitted to NeurIPS 2026 must be prepared according to the % instructions presented here. Papers may only be up to {\bf nine} pages long, including figures. \textbf{Papers that exceed the page limit will not be reviewed (or in any other way considered) for presentation at the conference.} % Additional pages \emph{containing acknowledgments, references, checklist, and optional technical appendices} do not count as content pages.    % The margins in 2026 are the same as those in previous years.   % Authors are required to use the NeurIPS \LaTeX{} style files obtainable at the NeurIPS website as indicated below. Please make sure you use the current files and not previous versions. Tweaking the style files may be grounds for desk rejection.   % \subsection{Retrieval of style files}   % The style files for NeurIPS and other conference information are available on the website at % \begin{center} %   \url{https://neurips.cc}. % \end{center} % % The file \verb+neurips_2026.pdf+ contains these instructions and illustrates the various formatting requirements your NeurIPS paper must satisfy.   % The only supported style file for NeurIPS 2026 is \verb+neurips_2026.sty+, rewritten for \LaTeXe{}. \textbf{Previous style files for \LaTeX{} 2.09,  Microsoft Word, and RTF are no longer supported.}   % The \LaTeX{} style file contains three optional arguments:  % \begin{itemize} % \item \verb+final+, which creates a camera-ready copy,  % \item \verb+preprint+, which creates a preprint for submission to, e.g., arXiv, \item \verb+nonatbib+, which will not load the \verb+natbib+ package for you in case of package clash. % \end{itemize}   % \paragraph{Preprint option} % If you wish to post a preprint of your work online, e.g., on arXiv, using the NeurIPS style, please use the \verb+preprint+ option. This will create a nonanonymized version of your work with the text ``Preprint. Work in progress.'' in the footer. This version may be distributed as you see fit, as long as you do not say which conference it was submitted to. Please \textbf{do not} use the \verb+final+ option, which should \textbf{only} be used for papers accepted to NeurIPS.   % At submission time, please omit the \verb+final+ and \verb+preprint+ options. This will anonymize your submission and add line numbers to aid review. Please do \emph{not} refer to these line numbers in your paper as they will be removed during generation of camera-ready copies.   % The file \verb+neurips_2026.tex+ may be used as a ``shell'' for writing your paper. All you have to do is replace the author, title, abstract, and text of the paper with your own.   % The formatting instructions contained in these style files are summarized in Sections \ref{gen_inst}, \ref{headings}, and \ref{others} below.   % \section{General formatting instructions} % \label{gen_inst}   % The text must be confined within a rectangle 5.5~inches (33~picas) wide and % 9~inches (54~picas) long. The left margin is 1.5~inch (9~picas).  Use 10~point % type with a vertical spacing (leading) of 11~points.  Times New Roman is the % preferred typeface throughout, and will be selected for you by default. % Paragraphs are separated by \nicefrac{1}{2}~line space (5.5 points), with no % indentation.   % The paper title should be 17~point, initial caps/lower case, bold, centered % between two horizontal rules. The top rule should be 4~points thick and the % bottom rule should be 1~point thick. Allow \nicefrac{1}{4}~inch space above and % below the title to rules. All pages should start at 1~inch (6~picas) from the % top of the page.   % For the final version, authors' names are set in boldface, and each name is % centered above the corresponding address. The lead author's name is to be listed % first (left-most), and the co-authors' names (if different address) are set to % follow. If there is only one co-author, list both author and co-author side by % side.   % Please pay special attention to the instructions in Section \ref{others} % regarding figures, tables, acknowledgments, and references.  % \section{Headings: first level} % \label{headings}   % All headings should be lower case (except for first word and proper nouns), % flush left, and bold.   % First-level headings should be in 12-point type.   % \subsection{Headings: second level}   % Second-level headings should be in 10-point type.   % \subsubsection{Headings: third level}   % Third-level headings should be in 10-point type.   % \paragraph{Paragraphs}   % There is also a \verb+\paragraph+ command available, which sets the heading in % bold, flush left, and inline with the text, with the heading followed by 1\,em % of space.   % \section{Citations, figures, tables, references} % \label{others}   % These instructions apply to everyone.   % \subsection{Citations within the text}   % The \verb+natbib+ package will be loaded for you by default.  Citations may be % author/year or numeric, as long as you maintain internal consistency.  As to the % format of the references themselves, any style is acceptable as long as it is % used consistently.   % The documentation for \verb+natbib+ may be found at % \begin{center} %   \url{http://mirrors.ctan.org/macros/latex/contrib/natbib/natnotes.pdf} % \end{center} % Of note is the command \verb+\citet+, which produces citations appropriate for % use in inline text.  For example, % \begin{verbatim} %    \citet{hasselmo} investigated\dots % \end{verbatim} % produces % \begin{quote} %   Hasselmo, et al.\ (1995) investigated\dots % \end{quote}   % If you wish to load the \verb+natbib+ package with options, you may add the % following before loading the \verb+neurips_2026+ package: % \begin{verbatim} %    \PassOptionsToPackage{options}{natbib} % \end{verbatim}   % If \verb+natbib+ clashes with another package you load, you can add the optional % argument \verb+nonatbib+ when loading the style file: % \begin{verbatim} %    \usepackage[nonatbib]{neurips_2026} % \end{verbatim}   % As submission is double blind, refer to your own published work in the third % person. That is, use ``In the previous work of Jones et al.\ [4],'' not ``In our % previous work [4].'' If you cite your other papers that are not widely available % (e.g., a journal paper under review), use anonymous author names in the % citation, e.g., an author of the form ``A.\ Anonymous'' and include a copy of the anonymized paper in the supplementary material.   % \subsection{Footnotes}   % Footnotes should be used sparingly.  If you do require a footnote, indicate % footnotes with a number\footnote{Sample of the first footnote.} in the % text. Place the footnotes at the bottom of the page on which they appear. % Precede the footnote with a horizontal rule of 2~inches (12~picas).   % Note that footnotes are properly typeset \emph{after} punctuation % marks.\footnote{As in this example.}   % \subsection{Figures}   % \begin{figure} %   \centering %   \fbox{\rule[-.5cm]{0cm}{4cm} \rule[-.5cm]{4cm}{0cm}} %   \caption{Sample figure caption. Explain what the figure shows and add a key take-away message to the caption.} % \end{figure}   % All artwork must be neat, clean, and legible. Lines should be dark enough for %  reproduction purposes. The figure number and caption always appear after the % figure. Place one line space before the figure caption and one line space after % the figure. The figure caption should be lower case (except for the first word and proper nouns); figures are numbered consecutively.   % You may use color figures.  However, it is best for the figure captions and the % paper body to be legible if the paper is printed in either black/white or in % color.   % \subsection{Tables}   % All tables must be centered, neat, clean, and legible.  The table number and % title always appear before the table.  See Table~\ref{sample-table}.   % Place one line space before the table title, one line space after the % table title, and one line space after the table. The table title must % be lower case (except for the first word and proper nouns); tables are % numbered consecutively.   % Note that publication-quality tables \emph{do not contain vertical rules}. We % strongly suggest the use of the \verb+booktabs+ package, which allows for % typesetting high-quality, professional tables: % \begin{center} %   \url{https://www.ctan.org/pkg/booktabs} % \end{center} % This package was used to typeset Table~\ref{sample-table}.   % \begin{table} %   \caption{Sample table caption. Explain what the table shows and add a key take-away message to the caption.} %   \label{sample-table} %   \centering %   \begin{tabular}{lll} %     \toprule %     \multicolumn{2}{c}{Part}                   \\ %     \cmidrule(r){1-2} %     Name     & Description     & Size ($\mu$m) \\ %     \midrule %     Dendrite & Input terminal  & $\approx$100     \\ %     Axon     & Output terminal & $\approx$10      \\ %     Soma     & Cell body       & up to $10^6$  \\ %     \bottomrule %   \end{tabular} % \end{table}  % \subsection{Math} % Note that display math in bare TeX commands will not create correct line numbers for submission. Please use LaTeX (or AMSTeX) commands for unnumbered display math. (You really shouldn't be using \$\$ anyway; see \url{https://tex.stackexchange.com/questions/503/why-is-preferable-to} and \url{https://tex.stackexchange.com/questions/40492/what-are-the-differences-between-align-equation-and-displaymath} for more information.)  % \subsection{Final instructions}  % Do not change any aspects of the formatting parameters in the style files.  In % particular, do not modify the width or length of the rectangle the text should % fit into, and do not change font sizes. Please note that pages should be % numbered.   % \section{Preparing PDF files}   % Please prepare submission files with paper size ``US Letter,'' and not, for % example, ``A4.''   % Fonts were the main cause of problems in the past years. Your PDF file must only % contain Type 1 or Embedded TrueType fonts. Here are a few instructions to % achieve this.   % \begin{itemize}   % \item You should directly generate PDF files using \verb+pdflatex+.   % \item You can check which fonts a PDF files uses.  In Acrobat Reader, select the %   menu Files$>$Document Properties$>$Fonts and select Show All Fonts. You can %   also use the program \verb+pdffonts+ which comes with \verb+xpdf+ and is %   available out-of-the-box on most Linux machines.   % \item \verb+xfig+ ``patterned'' shapes are implemented with bitmap fonts.  Use %   "solid" shapes instead.   % \item The \verb+\bbold+ package almost always uses bitmap fonts.  You should use %   the equivalent AMS Fonts: % \begin{verbatim} %    \usepackage{amsfonts} % \end{verbatim} % followed by, e.g., \verb+\mathbb{R}+, \verb+\mathbb{N}+, or \verb+\mathbb{C}+ % for $\mathbb{R}$, $\mathbb{N}$ or $\mathbb{C}$.  You can also use the following % workaround for reals, natural and complex: % \begin{verbatim} %    \newcommand{\RR}{I\!\!R} %real numbers %    \newcommand{\Nat}{I\!\!N} %natural numbers %    \newcommand{\CC}{I\!\!\!\!C} %complex numbers % \end{verbatim} % Note that \verb+amsfonts+ is automatically loaded by the \verb+amssymb+ package.   % \end{itemize}   % If your file contains type 3 fonts or non embedded TrueType fonts, we will ask % you to fix it.   % \subsection{Margins in \LaTeX{}}   % Most of the margin problems come from figures positioned by hand using % \verb+\special+ or other commands. We suggest using the command % \verb+\includegraphics+ from the \verb+graphicx+ package. Always specify the % figure width as a multiple of the line width as in the example below: % \begin{verbatim} %    \usepackage[pdftex]{graphicx} ... %    \includegraphics[width=0.8\linewidth]{myfile.pdf} % \end{verbatim} % See Section 4.4 in the graphics bundle documentation % (\url{http://mirrors.ctan.org/macros/latex/required/graphics/grfguide.pdf})   % A number of width problems arise when \LaTeX{} cannot properly hyphenate a % line. Please give LaTeX hyphenation hints using the \verb+\-+ command when % necessary.  % \begin{ack} % Use unnumbered first level headings for the acknowledgments. All acknowledgments % go at the end of the paper before the list of references. Moreover, you are required to declare % funding (financial activities supporting the submitted work) and competing interests (related financial activities outside the submitted work). % More information about this disclosure can be found at: \url{https://neurips.cc/Conferences/2026/PaperInformation/FundingDisclosure}.   % Do {\bf not} include this section in the anonymized submission, only in the final paper. You can use the \texttt{ack} environment provided in the style file to automatically hide this section in the anonymized submission. % \end{ack}  % \section*{References}  \bibliography{neurips2026/references} \bibliographystyle{plain}   % References follow the acknowledgments in the camera-ready paper. Use unnumbered first-level heading for % the references. Any choice of citation style is acceptable as long as you are % consistent. It is permissible to reduce the font size to \verb+small+ (9 point) % when listing the references. % Note that the Reference section does not count towards the page limit. % \medskip   % { % \small   % [1] Alexander, J.A.\ \& Mozer, M.C.\ (1995) Template-based algorithms for % connectionist rule extraction. In G.\ Tesauro, D.S.\ Touretzky and T.K.\ Leen % (eds.), {\it Advances in Neural Information Processing Systems 7}, % pp.\ 609--616. Cambridge, MA: MIT Press.   % [2] Bower, J.M.\ \& Beeman, D.\ (1995) {\it The Book of GENESIS: Exploring %   Realistic Neural Models with the GEneral NEural SImulation System.}  New York: % TELOS/Springer--Verlag.   % [3] Hasselmo, M.E., Schnell, E.\ \& Barkai, E.\ (1995) Dynamics of learning and % recall at excitatory recurrent synapses and cholinergic modulation in rat % hippocampal region CA3. {\it Journal of Neuroscience} {\bf 15}(7):5249-5262. % }   %%%%%%%%%%%%%%%%%%%%%%%%%%%%%%%%%%%%%%%%%%%%%%%%%%%%%%%%%%%%  \appendix  \section{Quick-eval Solver Comparison Results}  \begin{table}[h] \centering \captionsetup{font=small} \caption{Preliminary Solver comparison results. \textcolor{red}{*PAR-2 on selected problems.} % to further speedup the preliminary study, thus may not fully correlate with full evaluation in Table~\ref{tab:solver_comparison}.} } % \vspace{-1.0em} \resizebox{.99\textwidth}{!}{ \begin{tabular}{|l|l|c|c|} \hline \textbf{Solver} & \textbf{Heuristics} & \textbf{Best PAR-2* $\downarrow$} & \textbf{Base PAR-2* $\downarrow$} \\ \hline \multirow{3}{*}{Kissat\_vsa}   & init\_activityref                          & 1066.66 & \multirow{3}{*}{575.03} \\   & rank\_vivification\_candidates\_by\_activity & 763.8   & \\   & sort\_vivification\_candidates\_by\_activity & 760.19  & \\ \hline Kissat\_pred     & predict                & 893.65 & 654.98 \\ \hline Kissat-Corephase & reset\_phases\_mab part2 & 603.99 & 513.91 \\ \hline \multirow{3}{*}{AE\_kissat2025\_MAB}   & kissat\_restarting             & 629.15 & \multirow{3}{*}{528.76} \\   & \textbf{restart\_mab}                   & \textbf{518.09} & \\   & \textbf{kissat\_bump\_score\_increment}  & \textbf{475.58} & \\ \hline \multirow{3}{*}{Kissat-Cure}   & \textbf{kissat\_dynamicsat}     & \textbf{469.44} & \multirow{3}{*}{531.05} \\   % & kissat\_decide\_phase  & 522.5  & \\   & kissat\_reboot\_direct & 566.63 & \\   & kissat\_coldrestarting & 678.23 & \\ \hline \multirow{5}{*}{Kissat\_Corephase\_Coreward}   & walking\_phase                       & 874.72 & \multirow{5}{*}{740.72} \\   & \textbf{rephase\_conflict}                    & \textbf{603.94} & \\   % & largest\_score\_unassigned\_variable  & 602.31 & \\   % & kissat\_analyze                      & 583.21 & \\   & rephase\_walking                     & 708.82 & \\   & kissat\_rephase                      & 738.37 & \\   & reset\_phases                        & 739.68 & \\ \hline \end{tabular} } % \caption{Solver performance comparison} \label{tab:solver_comparison_quick} \end{table}  \section{Implementation Details} \label{app:implementation}  \subsection{Structured Step Decomposition and Proof Validation} \label{app:step_decomposition} The algorithm is constrained to an ordered sequence of sub-steps demarcated by ``Step:'' markers. Their step decompositions then serve as the basis for the next round of structured local search.   After evaluation, our pipeline runs the DRAT-trim~\cite{heule2016drat} proof checking utility on the top-performing algorithms to certify the proofs of their UNSAT results. We only advance a solver through the pipeline if, after one loop of mutant generation, it is both the top-performing solver in its team and produces valid results on all benchmark instances.   \subsection{Coder Scaffolding and Signal Reference Construction} \label{app:signal_reference} The scaffolding of the coder model has two main requirements: the implementation must be \textit{semantically faithful} to the natural-language specification, and it must respect the affordances of the target function's role within Kissat. Specifically, it must work within the accessible state, return type, and surrounding control flow of the target function. Naively prompting a strong coder model satisfies neither of these requirements reliably.  The reference was constructed through a hybrid manual-LLM procedure. We first ran the pipeline without any reference and logged the recurring solver-state requirements that natural-language algorithms assumed, including current decision level, variable phase, Kissat macros, and similar constructs. This seed list was then passed, together with the solver source code, to an LLM tasked with extending it to cover edge cases the initial sample had missed. We iterated this loop until the compilation rate in each team reached 90\%, which in practice required only a single refinement step. Our reference is deliberately minimal to protect against biasing the coder model toward implementations that lean heavily on the signals that are included.  \subsection{LLM-Generated Analysis for Memory Bank Records} \label{app:analysis_generation} Crucially, the same step-aware representation used during mutant generation serves as the key for persistent memory bank records, so retrieval at mutation time is guaranteed to be semantically aligned with the query that produced the record. Analyses are produced offline by a separate LLM call once an iteration's evaluation has finished. For each selected parent-mutant pair, the analyzer receives the parent description, the mutated step, the mutant description, both PAR-2 scores, the relative change, and a unified code diff between the parent and child solver implementations, and is asked to explain the improvement or degradation in concrete, code-grounded terms. Improvements and regressions use separate prompt templates so the analyzer focuses on the right outcome; pairs are batched in a single request so the LLM can identify relevant patterns across them, and only responses that pass schema validation are written to the index. Each persisted record therefore captures not only what happened but also why it happened, in language directly consumable by the downstream mutation prompt.  \subsection{Stratified Benchmark Subset Construction} \label{app:benchmark_subset} To enable rapid iteration during heuristic search, we constructed a stratified 50-instance subset of the SAT Competition 2025 benchmark~\cite{codel2025sat}. We first ran the unmodified base solver on the full 400-instance set with a 5000s wall-clock timeout to obtain a per-instance solving time, then bucketed instances by runtime into four difficulty tiers: Easy ($<$10s), Medium (10--1000s), Hard (1000--5000s), and Timeout. We sampled 50 instances proportionally using a fixed random seed (42) for reproducibility. The resulting subset preserves the full benchmark's difficulty profile while substantially reducing per-iteration evaluation cost. The code is available in our submission.  \subsection{Target Functions}  \begin{enumerate}[leftmargin=*] \item Only modified or new functions. \item Functions must implement the core algorithm-level innovations. \item We exclude deterministic functions whose specification fully determines their behavior, leaving no design space to explore. \end{enumerate}     The table below details the solvers and heuristics, including their design rationales and respective roles. \begin{table}[h] \centering \captionsetup{font=small} \caption{Experimental targets: large-scale SAT solver variants and their core heuristic functions selected for evolution. \textcolor{red}{Todo: we may use LLM to identify these heuristics, beyond the scope of this work. Even we manually picked these heuristics, zero-shot transfer KissatEvolve to different heuristics is still non-trivial.}} \resizebox{.99\textwidth}{!}{ \begin{tabular}{|l|l|l|} \hline \textbf{Solver} & \textbf{Targeted Heuristic Functions} & \textbf{Focus Area} \\ \hline \multirow{3}{*}{Kissat\_vsa~\cite{li2025kissatvsa}}   & init\_activityref                            & \multirow{3}{*}{Variable activity \& vivification} \\   & rank\_vivification\_candidates\_by\_activity & \\   & sort\_vivification\_candidates\_by\_activity & \\ \hline \multirow{4}{*}{Kissat\_pred~\cite{luo2025kissatpred}}   & kissat\_set\_configuration  & \multirow{4}{*}{Predictive \& configuration parameters} \\   & predict                     & \\   & extract\_features           & \\   & extract\_sequence\_features & \\ \hline \multirow{8}{*}{\shortstack[l]{Kissat-Corephase~\cite{chen2025kissatcorephase} \\ \& Kissat\_Corephase\_Coreward}}   & reset\_phases\_mab                    & \multirow{8}{*}{Phase-saving \& walking behaviors} \\   & walking\_phase                        & \\   & rephase\_conflict                     & \\   & largest\_score\_unassigned\_variable  & \\   & kissat\_analyze                       & \\   & rephase\_walking                      & \\   & reset\_phases                         & \\   & kissat\_rephase                       & \\ \hline \multirow{4}{*}{Kissat-Cure~\cite{cai2025kissatcure}}   & kissat\_dynamicsat       & \multirow{4}{*}{Dynamic phase decisions \& restarts} \\   & kissat\_decide\_phase    & \\   & kissat\_reboot\_direct   & \\   & kissat\_coldrestarting   & \\ \hline \multirow{4}{*}{AE\_kissat2025\_MAB}   & kissat\_restarting               & \multirow{4}{*}{Restarting \& scoring mechanics} \\   & restart\_mab                     & \\   & kissat\_bump\_score\_increment   & \\   & kissat\_rescale\_scores          & \\ \hline \end{tabular} } \label{tab:experimental_targets} \end{table}    \section{Cross-over Stage Analysis} \label{app:cross_over_detail} The cross-over stage utilizes parents from the mutation stage as inputs and combines complementary strengths across algorithms from different mutation groups. For each parent, the cross-over LLM first performs causal analysis to identify its strengths, weaknesses, and key mechanisms. Then, the LLM proposes the most promising pairings based on the causal reports and past combination outcomes retrieved from the memory bank, with only high-scoring proposals proceeding to crossover. For each accepted pair, the LLM generates a new offspring algorithm in a step-by-step structured format, which is then translated into C code and evaluated. Offspring that fail to improve over the best parent PAR-2 are discarded, ensuring the population only retains beneficial combinations.  To investigate why cross-over underperforms, we examine two heuristic functions under the AE\_kissat2025\_MAB solver, kissat\_restarting and kissat\_bump\_score\_increment, which individually improved over their baselines during the mutation stage. Adding the cross-over stage degraded performance by +83.83 and +62.14 PAR-2 respectively. We attribute this to two factors. First, reconciling two independently evolved implementations in a single code-generation call is a harder task than translating a single mutant. Second, cross-over's diversity advantage is mainly achieved by combining parents from different lineages. This is substantially weakened by our diversified mutation design. With temperature annealing, cosine similarity monitoring, and step-level controlled mutations over multiple iterations, the mutation stage already covers a broad region of the heuristic space, making cross-over's  contribution limited while its cost remains high.  \section{Prompts} \subsection{Leader Prompt} \label{app:leader_prompt} % --- Leader Prompt --- \begin{promptbox}[title={Prompt: Leader (Algorithm Generation)}] \ttfamily\small You are designing a novel and diverse heuristic strategy on top of the Kissat solver which will optimize performance on the SAT competition benchmarks (PAR-2 scoring).  This algorithm (conceptual, not code) should modify the target function below. Your algorithm should demonstrate the same level of innovation and diversity as past SAT competition winners.  Be specific about the mechanism and logic flow. Include concrete thresholds and parameters where necessary, but prioritize describing a structurally unique decision process over tuning constants.  Your algorithm does not need to be long, just intelligent, innovative, and mechanistically distinct from standard approaches.  \medskip \textbf{Target Function}  \texttt{FUNCTION\_NAME\_PLACEHOLDER}  \medskip \textbf{Output Format (JSON only)}  \begin{verbatim} {   "name": "Brief algorithm name (<=6 words)",   "algorithm": "Complete algorithmic description. Step 1: ... Step 2: ...",   "reason": "1-2 brief lines about why this improves performance.",   "target_function": "exact_function_name_from_list_above" } \end{verbatim}  \textit{IMPORTANT: \texttt{target\_function} must exactly match the target function name listed above (no parens).}\\ \textit{IMPORTANT: \texttt{algorithm} must specifically split substeps by Step markers ``Step 1: \ldots\ Step 2: \ldots''}\\ \textit{Your algorithm will be implemented by a coder model. Make sure they generally have the requisite information in your algorithmic description to be able to implement it faithfully.} \end{promptbox}  \subsection{Mutant Prompt} \label{app:mutant_prompt} \begin{promptbox}[title={Prompt: Mutant (Algorithm Generation)}] \ttfamily\small You are designing a variant heuristic strategy on top of the Kissat solver which will optimize performance on the SAT competition benchmarks (PAR-2 scoring).  \medskip Here is the given strategy that you are creating a variant of:  \texttt{\{leader\_algorithm\}}  \medskip You will be targeting a specific step within the leader algorithm. Your step will be:  \texttt{\{target\_step\_num\}}  \medskip \texttt{\{experience\_pool\_section\}}  \medskip Generate a variant of this algorithm that makes changes to this specific step. Your changes can make it longer or shorter depending on your expert SAT solving opinion. Include in your reason the step of the original algorithm.  \medskip \textbf{Required Output JSON}  \begin{verbatim} {   "name": "Brief algorithm name (<=6 words)",   "algorithm": "Complete algorithmic description.",   "reason": "Why this variant improves on the original." } \end{verbatim}  \textit{IMPORTANT: Your algorithm should be the full algorithm with your step change replaced, not just the specific step.} \end{promptbox}   \subsection{Code Generation Prompt} \label{app:code_gen_prompt} % --- Coder Prompt --- \begin{promptbox}[title={Prompt: Coder (C Implementation)}] \ttfamily\small You are an expert in writing C code for Kissat-based SAT Solvers and your task is to faithfully implement the given algorithm below so that it compiles and implements the strategy that the algorithm describes for the improvement of the SAT Solver.  \medskip \textbf{Requirements} \begin{itemize}   \item Output complete function with signature and body   \item Must compile with existing Kissat codebase   \item \textbf{CRITICAL}: Do NOT invent or hallucinate solver fields   \item \textbf{CRITICAL}: Only the function body is injected into the source file. You CANNOT add new \texttt{\#include} directives, file-scope declarations, or helper functions. All persistent state must use \texttt{static} local variables inside your function. \end{itemize}  \medskip \textbf{Kissat API Reference}  \textit{1. Solver State Access (\texttt{kissat *solver})}  \medskip \textit{Search State:}  \begin{tabular}{lll} \toprule Field & Type & Description \\ \midrule \texttt{solver->level} & \texttt{unsigned} & Current decision level (0 = root) \\ \texttt{solver->stable} & \texttt{bool} & true=stable, false=focused \\ \texttt{solver->vars} & \texttt{unsigned} & Total number of variables \\ \texttt{solver->unassigned} & \texttt{unsigned} & Count of unassigned variables \\ \texttt{solver->ticks} & \texttt{uint64\_t} & Overall work counter \\ \texttt{solver->warming} & \texttt{bool} & True during initial local-search warming \\ \bottomrule \end{tabular}  \medskip \textit{Phase Arrays:}  \begin{tabular}{lll} \toprule Field & Type & Description \\ \midrule \texttt{solver->phases.saved[idx]} & \texttt{value} & Saved phase (-1, 0, or 1). 0 = uninitialized \\ \texttt{solver->phases.target[idx]} & \texttt{value} & Target phase \\ \texttt{solver->phases.best[idx]} & \texttt{value} & Best phase (from best overall assignment) \\ \bottomrule \end{tabular}  \smallskip The \texttt{value} type is \texttt{typedef signed char value} --- it holds -1, 0, or 1.  \medskip \textit{Assignment State:}  \begin{tabular}{lll} \toprule Field & Type & Description \\ \midrule \texttt{solver->values} & \texttt{value *} & Truth values indexed by literal. \\ \texttt{solver->assigned} & \texttt{assigned *} & Per-variable assignment info \\ \bottomrule \end{tabular}  \smallskip The \texttt{assigned} struct fields: \texttt{.level} (decision level), \texttt{.trail} (position), \texttt{.reason} (clause ref or \texttt{DECISION\_REASON}).  \medskip \textit{Scoring Fields:}  \begin{tabular}{lll} \toprule Field & Type & Description \\ \midrule \texttt{solver->scores} & \texttt{heap} & VSIDS score heap \\ \texttt{solver->scinc} & \texttt{double} & Current VSIDS score increment \\ \bottomrule \end{tabular}  \medskip \textit{Multi-Armed Bandit (MAB) Fields:}  \begin{tabular}{lll} \toprule Field & Type & Description \\ \midrule \texttt{solver->mab} & \texttt{bool} & MAB enabled flag \\ \texttt{solver->heuristic} & \texttt{unsigned} & Current heuristic (0=VSIDS, 1=CHB) \\ \bottomrule \end{tabular}  \medskip \textit{Random Number Generator:}  \begin{tabular}{lll} \toprule Field & Type & Description \\ \midrule \texttt{solver->random} & \texttt{generator} (\texttt{uint64\_t}) & RNG state \\ \bottomrule \end{tabular}  \medskip \textit{2. Statistics (\texttt{solver->statistics})}  \medskip Access macros: \texttt{CONFLICTS}, \texttt{DECISIONS}, \texttt{INC(name)}, \texttt{ADD(name, n)}, \texttt{GET(name)}.  \begin{tabular}{ll} \toprule Counter & Description \\ \midrule \texttt{conflicts} & Total conflicts \\ \texttt{decisions} & Total decisions \\ \texttt{propagations} & Total unit propagations \\ \texttt{restarts} & Total restarts \\ \texttt{clauses\_learned} & Learned clauses added \\ \texttt{clauses\_irredundant} & Original clauses \\ \texttt{clauses\_redundant} & Redundant learned clauses \\ \texttt{target\_decisions} & Decisions using target phase \\ \texttt{saved\_decisions} & Decisions using saved phase \\ \texttt{initial\_decisions} & Decisions using initial/default phase \\ \texttt{switched} & Mode switch counter \\ \bottomrule \end{tabular}  \smallskip Glue distribution: \texttt{solver->statistics.used[mode].glue[lbd\_value]} (\texttt{uint64\_t}); mode: 0=focused, 1=stable; lbd\_value: 0--127.  \medskip \textit{3. Averages (Exponential Moving Averages)}  \texttt{AVERAGE(fast\_glue)}, \texttt{AVERAGE(slow\_glue)}, \texttt{AVERAGE(level)}, \texttt{AVERAGE(trail)}, \texttt{AVERAGE(decision\_rate)}.  \medskip \textit{4. Options --- \texttt{GET\_OPTION(name)}}  \begin{tabular}{llll} \toprule Option & Default & Range & Description \\ \midrule \texttt{forcephase} & 0 & 0--1 & Force initial phase \\ \texttt{phase} & 1 & 0--1 & Initial phase polarity \\ \texttt{phasesaving} & 1 & 0--1 & Enable phase saving \\ \texttt{target} & 1 & 0--2 & Use target phases \\ \texttt{decay} & 50 & 1--200 & Per mille scores decay \\ \texttt{tier1} & 2 & 1--100 & Tier 1 glue limit \\ \texttt{tier2} & 6 & 1--1000 & Tier 2 glue limit \\ \texttt{heuristic} & 0 & 0--1 & Scoring heuristic (0=VSIDS, 1=CHB) \\ \texttt{restart} & 1 & 0--1 & Enable restarts \\ \texttt{restartint} & 50 & 1--10000 & Base restart interval \\ \texttt{restartmargin} & 10 & 0--25 & Fast/slow glue margin (percent) \\ \bottomrule \end{tabular}  \medskip \textit{5. Clause Structure}  \begin{verbatim} struct clause {   unsigned glue : 19;   // LBD value   bool garbage : 1;     // Marked for deletion   bool reason : 1;      // Used as reason clause   bool redundant : 1;   // Learned clause   bool shrunken : 1;    // Clause was minimized   unsigned used : 5;    // Usage count (0-31)   unsigned searched;    // Search counter   unsigned size;        // Number of literals   unsigned lits[3];     // Literal array }; \end{verbatim}  Iteration: \begin{verbatim} for (all_clauses(c)) { ... } for (all_literals_in_clause(lit, c)) { ... } \end{verbatim}  \medskip \textit{6. Available Functions}  \begin{tabular}{ll} \toprule Function & Description \\ \midrule \texttt{kissat\_pick\_random(\&solver->random, low, high)} & Returns \texttt{unsigned} in [low, high) \\ \texttt{kissat\_pick\_double(\&solver->random)} & Returns \texttt{double} in [0, 1) \\ \texttt{kissat\_pick\_bool(\&solver->random)} & Returns random \texttt{bool} \\ \texttt{kissat\_next\_random32(\&solver->random)} & Returns \texttt{unsigned} (32-bit) \\ \texttt{kissat\_get\_heap\_score(heap, idx)} & Returns \texttt{double} score for variable \\ \texttt{kissat\_max\_heap(heap)} & Returns \texttt{unsigned} idx of max var \\ \texttt{kissat\_logn(count)} & Returns \texttt{double}: $\log_n(\text{count})$ \\ \texttt{kissat\_sqrt(count)} & Returns \texttt{double}: $\sqrt{\text{count}}$ \\ \texttt{kissat\_nlogpown(count, p)} & Returns \texttt{double}: $n \cdot (\log_n)^p$ \\ \texttt{kissat\_calloc(solver, n, size)} & Allocate $n \times \text{size}$ zeroed bytes \\ \texttt{kissat\_malloc(solver, bytes)} & Allocate bytes \\ \texttt{kissat\_realloc(solver, ptr, old, new)} & Reallocate \\ \texttt{kissat\_free(solver, ptr, bytes)} & Free memory \\ \bottomrule \end{tabular}  \medskip \textit{7. Utility Macros}  \begin{tabular}{ll} \toprule Macro & Description \\ \midrule \texttt{IDX(lit)} & Extract variable index from literal \\ \texttt{LIT(idx)} & Create positive literal from variable index \\ \texttt{NOT(lit)} & Negate literal \\ \texttt{VALUE(lit)} & Current truth value (-1, 0, 1) --- takes a literal \\ \texttt{ACTIVE(idx)} & True if variable is active \\ \texttt{VARS} & \texttt{solver->vars} \\ \texttt{LITS} & \texttt{2 * solver->vars} \\ \texttt{INITIAL\_PHASE} & \texttt{(GET\_OPTION(phase) ? 1 : -1)} \\ \texttt{FRAME(level)} & Get frame at decision level \\ \texttt{SCORES} & \texttt{\&solver->scores} \\ \bottomrule \end{tabular}  \medskip \textit{8. Persistent State Pattern}  \begin{verbatim} int kissat_decide_phase(kissat *solver, unsigned idx) {     static double *my_array = 0;     static unsigned my_array_size = 0;     if (!my_array || VARS > my_array_size) {         unsigned new_size = VARS + 1000;         double *new_arr = kissat_calloc(solver, new_size, sizeof(double));         if (my_array) {             for (unsigned i = 0; i < my_array_size; i++)                 new_arr[i] = my_array[i];             kissat_free(solver, my_array, my_array_size * sizeof(double));         }         my_array = new_arr;         my_array_size = new_size;     } } \end{verbatim}  \medskip \textit{9. Critical Rules --- What NOT to do} \begin{itemize}   \item \textbf{DO NOT} invent solver fields --- only use fields documented above   \item \textbf{DO NOT} create \texttt{kissat\_*} prefixed variables (reserved for API). Name static variables without the \texttt{kissat\_} prefix (e.g., \texttt{phase\_scores}, \texttt{flip\_count})   \item \textbf{DO NOT} use \texttt{math.h} functions (\texttt{fabs}, \texttt{sqrt}, \texttt{log}, \texttt{log2}) --- use \texttt{(x > 0 ? x : -x)} or \texttt{kissat\_logn()}/\texttt{kissat\_sqrt()} instead   \item \textbf{DO NOT} use \texttt{stdlib.h}/\texttt{string.h}/\texttt{stdio.h} --- use \texttt{kissat\_malloc}/\texttt{calloc}/\texttt{free}   \item \textbf{NO} \texttt{solver->lbd} --- LBD is per-clause as \texttt{clause->glue}   \item \textbf{NO} \texttt{solver->conflict} --- use \texttt{CONFLICTS} macro   \item \textbf{NO} \texttt{solver->clauses.size} --- iterate with \texttt{for (all\_clauses(c))} \end{itemize}  Return value must be \texttt{1} (positive phase) or \texttt{-1} (negative phase). Never return \texttt{0}.  \medskip \textbf{Algorithm to Implement}  \texttt{ALGORITHM\_PLACEHOLDER}  \medskip \textbf{Output Format}  Wrap your complete function in \texttt{<code>} tags (no other text outside tags):  \begin{verbatim} <code> void FUNCTION_NAME_PLACEHOLDER(kissat *solver) {     // Your implementation here } </code> \end{verbatim}  \textit{IMPORTANT: Wrap your code between \texttt{<code>} \texttt{</code>} tags so that it can be properly parsed.}  \medskip \textbf{Baseline Reference}  The full source of \texttt{FILENAME\_PLACEHOLDER} is provided. Your generated \texttt{FUNCTION\_NAME\_PLACEHOLDER} function will replace the existing one; ensure it compiles with the rest of the code and faithfully implements the algorithm.  \texttt{SOURCE\_CODE\_PLACEHOLDER} \end{promptbox}       \section{Examples of Generated Heuristics} \label{app:examples} \subsection{Example 1: Kissat\_Cure Solver, kissat\_dynamicsat Function}  \paragraph{Parent: Natural Language} Topological Locality Adaptive Mode Switching: Step 1: Maintain a 'Conflict Locality Index' (CLI) calculated every 1024 conflicts as the average literal-distance between variables in the conflict clause using the initial variable order as a 1D projection of the graph. Step 2: If the CLI is below the threshold of 0.25 (indicating high cluster locality), the solver enters 'Focused-Locality' mode, where the VSIDS decay factor is tightened to 0.75 and the restart interval is fixed to a rapid Luby sequence to intensify search within the detected cluster. Step 3: If CLI exceeds 0.6 (indicating global constraints), the solver switches to 'Diffuse-Global' mode, utilizing a CHB-style heuristic with a relaxed decay of 0.98 and phase-saving that favors 'Target' phases over 'Best' phases to encourage exploration. Step 4: Every 2048 conflicts, evaluate the 'LBD Slope' (derivative of the Literal Block Distance of learned clauses); if the slope is positive while in Focused-Locality mode, force a 'Phase Flip' on the top 5\% of variables to escape a potential local search trap.  \paragraph{Parent: Code Implementation} \begin{codebox} \begin{lstlisting}[style=cstyle] #include "internal.h" #include <math.h> #include <stdlib.h> #include <string.h>  void kissat_dynamicsat(kissat *solver) {     static uint64_t last_cli_count = 0;     static uint64_t last_lbd_count = 0;     static uint64_t last_conflicts = 0;     static double last_lbd_avg = 0;     static int current_mode = 0; // 0: Normal, 1: Focused-Locality, 2: Diffuse-Global      uint64_t conflicts = CONFLICTS;      // Step 1: CLI via span proxy every 1024 conflicts     if (conflicts >= last_cli_count + 1024) {         last_cli_count = conflicts;         double total_locality = 0;         uint64_t count = 0;         for (all_clauses(c)) {             if (c->garbage || !c->redundant) continue;             unsigned min_v = VARS, max_v = 0;             for (all_literals_in_clause(lit, c)) {                 unsigned v = IDX(lit);                 if (v < min_v) min_v = v;                 if (v > max_v) max_v = v;             }             if (max_v > min_v && c->size > 1) {                 total_locality += (double)(max_v - min_v) / (double)(c->size - 1);                 count++;             }         }         double cli = (count > 0) ? (total_locality / (double)count) : ((double)VARS * 0.5);         cli /= (double)VARS;          if (cli < 0.25) current_mode = 1;      // Focused-Locality         else if (cli > 0.6) current_mode = 2;  // Diffuse-Global         else current_mode = 0;                  // Normal     }      // Adjust scinc scaling based on mode     if (conflicts > last_conflicts) {         if (current_mode == 1) solver->scinc *= 1.05;         else if (current_mode == 2) solver->scinc *= 0.95;         last_conflicts = conflicts;     }      // Step 2: Focused-Locality Mode     if (current_mode == 1) {         solver->stable = false;     }     // Step 3: Diffuse-Global Mode     else if (current_mode == 2) {         solver->stable = true;         for (all_variables(idx)) {             if (ACTIVE(idx)) {                 value target = solver->phases.target[idx];                 if (target != 0) solver->phases.saved[idx] = target;             }         }     }      // Step 4: LBD Slope via histogram + sampling-based phase flip every 2048 conflicts     if (conflicts >= last_lbd_count + 2048) {         last_lbd_count = conflicts;         double lbd_sum = 0;         uint64_t lbd_n = 0;         for (int m = 0; m < 2; m++) {             for (int g = 0; g < 128; g++) {                 uint64_t val = solver->statistics.used[m].glue[g];                 lbd_sum += (double)g * val;                 lbd_n += val;             }         }         double avg_lbd = (lbd_n > 0) ? (lbd_sum / (double)lbd_n) : 0;         double slope = avg_lbd - last_lbd_avg;         last_lbd_avg = avg_lbd;          if (current_mode == 1 && slope > 0) {             double samples[100];             int sampled_count = 0;             for (int i = 0; i < 500 && sampled_count < 100; i++) {                 unsigned v = kissat_pick_random(&solver->random, 0, VARS);                 if (ACTIVE(v)) samples[sampled_count++] = kissat_get_heap_score(SCORES, v);             }             if (sampled_count > 10) {                 for (int i = 0; i < sampled_count; i++) {                     for (int j = i + 1; j < sampled_count; j++) {                         if (samples[i] < samples[j]) {                             double tmp = samples[i];                             samples[i] = samples[j];                             samples[j] = tmp;                         }                     }                 }                 double threshold = samples[sampled_count / 20]; // 95th percentile                 for (all_variables(idx)) {                     if (ACTIVE(idx) && kissat_get_heap_score(SCORES, idx) >= threshold) {                         if (solver->phases.saved[idx] == 0) solver->phases.saved[idx] = 1;                         solver->phases.saved[idx] *= -1;                     }                 }             }         }     } } \end{lstlisting} \end{codebox}  \paragraph{Mutant 1: Natural Language} Topological Locality with Target-Best Sync: Step 1: Maintain a 'Conflict Locality Index' (CLI) calculated every 1024 conflicts as the average literal-distance between variables in the conflict clause using the initial variable order as a 1D projection of the graph. Step 2: If the CLI is below the threshold of 0.25 (indicating high cluster locality), the solver enters 'Focused-Locality' mode, where the VSIDS decay factor is tightened to 0.75 and the restart interval is fixed to a rapid Luby sequence to intensify search within the detected cluster. Step 3: If CLI exceeds 0.6 (indicating global constraints), the solver switches to 'Diffuse-Global' mode, utilizing a CHB-style heuristic with a relaxed decay of 0.98 and phase-saving that favors 'Target' phases over 'Best' phases to encourage exploration. Step 4: Every 2048 conflicts, evaluate the 'LBD Slope' (derivative of the Literal Block Distance of learned clauses); if the slope is positive while in Focused-Locality mode, trigger a 'Target-Best Synchronization' where the current 'Target' phases are overwritten by 'Best' phases for variables with high connectivity (occurrence in the last 256 conflict clauses), and double the VSIDS increment for the next 512 conflicts to accelerate the discovery of a new local optimum.  \paragraph{Mutant 1: Code Implementation} \begin{codebox} \begin{lstlisting}[style=cstyle] #include "internal.h" #include <math.h> #include <stdlib.h> #include <string.h>  void kissat_dynamicsat(kissat *solver) {     static uint64_t last_cli_conf = 0;     static uint64_t last_lbd_conf = 0;     static double current_cli = 0.5;     static double last_avg_lbd = 0.0;     static int current_mode = 0; // 0: Normal, 1: Adaptive-Intense, 2: Diffuse-Global      uint64_t conflicts = CONFLICTS;      // Step 1: Maintain CLI every 1024 conflicts     if (conflicts >= last_cli_conf + 1024) {         double total_clause_locality = 0;         uint64_t redundant_count = 0;         for (all_clauses(c)) {             if (c->redundant && !c->garbage && c->size > 1) {                 double avg_idx = 0;                 for (all_literals_in_clause(lit, c)) {                     avg_idx += (double)IDX(lit);                 }                 avg_idx /= c->size;                  double var_dist = 0;                 for (all_literals_in_clause(lit_dist, c)) {                     double d = (double)IDX(lit_dist) - avg_idx;                     var_dist += (d > 0 ? d : -d);                 }                 total_clause_locality += (var_dist / c->size);                 redundant_count++;             }         }          if (redundant_count > 0 && VARS > 0) {             current_cli = (total_clause_locality / (double)redundant_count) / (double)VARS;         }         last_cli_conf = conflicts;          // Step 2: Adaptive-Intense Mode (CLI < 0.25)         if (current_cli < 0.25) {             current_mode = 1;             double multiplier = 0.95 - (0.25 - current_cli);             int decay_opt = (int)(1000.0 * (1.0 - multiplier));             if (decay_opt < 1) decay_opt = 1;             if (decay_opt > 200) decay_opt = 200;             kissat_set_option(solver, "decay", decay_opt);             kissat_set_option(solver, "restartmargin", 10);             kissat_set_option(solver, "stable", 1);         }         // Step 3: Diffuse-Global Mode (CLI > 0.6)         else if (current_cli > 0.6) {             current_mode = 2;             kissat_set_option(solver, "stable", 0);             kissat_set_option(solver, "decay", 20);             for (all_variables(idx)) {                 if (ACTIVE(idx)) {                     if (solver->phases.target[idx] != 0) {                         solver->phases.saved[idx] = solver->phases.target[idx];                     }                 }             }         } else {             current_mode = 0;         }     }      // Step 4: LBD Slope Evaluation every 2048 conflicts     if (conflicts >= last_lbd_conf + 2048) {         double current_avg_lbd = 0;         uint64_t red_lbd_count = 0;         for (all_clauses(c_lbd)) {             if (c_lbd->redundant && !c_lbd->garbage) {                 current_avg_lbd += c_lbd->glue;                 red_lbd_count++;             }         }         if (red_lbd_count > 0) current_avg_lbd /= red_lbd_count;         double lbd_slope = current_avg_lbd - last_avg_lbd;          if (lbd_slope > 0 && current_mode == 1) {             unsigned active_count = 0;             double max_score = 0;             for (all_variables(idx_score)) {                 if (ACTIVE(idx_score)) {                     active_count++;                     double s = kissat_get_heap_score(SCORES, idx_score);                     if (s > max_score) max_score = s;                 }             }              unsigned num_to_flip = active_count / 20; // top 5%             if (num_to_flip > 0) {                 double low = 0, high = max_score;                 for (int i = 0; i < 20; i++) {                     double mid = low + (high - low) / 2.0;                     unsigned count_above = 0;                     for (all_variables(idx_find)) {                         if (ACTIVE(idx_find) && kissat_get_heap_score(SCORES, idx_find) >= mid)                             count_above++;                     }                     if (count_above > num_to_flip) low = mid;                     else high = mid;                 }                 for (all_variables(idx_flip)) {                     if (ACTIVE(idx_flip) && kissat_get_heap_score(SCORES, idx_flip) >= low) {                         if (solver->phases.saved[idx_flip] == 0)                             solver->phases.saved[idx_flip] = 1;                         solver->phases.saved[idx_flip] *= -1;                     }                 }             }         }         last_avg_lbd = current_avg_lbd;         last_lbd_conf = conflicts;     } } \end{lstlisting} \end{codebox}  \paragraph{Mutant 2: Natural Language} Dynamic Decay and EMA-Triggered Locality: Step 1: Maintain a 'Conflict Locality Index' (CLI) calculated every 1024 conflicts as the average literal-distance between variables in the conflict clause using the initial variable order as a 1D projection of the graph. Step 2: If the CLI is below 0.25, the solver enters 'Adaptive-Intense' mode; here, the VSIDS decay is dynamically scaled as (0.95 - (0.25 - CLI)) to match the cluster density, and the restart policy is switched to an Exponential Moving Average (EMA) trigger with a threshold of 1.1 to force rapid, quality-sensitive restarts. Step 3: If CLI exceeds 0.6 (indicating global constraints), the solver switches to 'Diffuse-Global' mode, utilizing a CHB-style heuristic with a relaxed decay of 0.98 and phase-saving that favors 'Target' phases over 'Best' phases to encourage exploration. Step 4: Every 2048 conflicts, evaluate the 'LBD Slope' (derivative of the Literal Block Distance of learned clauses); if the slope is positive while in Focused-Locality mode, force a 'Phase Flip' on the top 5\% of variables to escape a potential local search trap.  \paragraph{Mutant 2: Code Implementation} \begin{codebox} \begin{lstlisting}[style=cstyle] #include "internal.h" #include <math.h> #include <stdlib.h> #include <string.h>  void kissat_dynamicsat(kissat *solver) {     static uint64_t last_cli_checkpoint = 0;     static uint64_t last_lbd_checkpoint = 0;     static double last_avg_lbd = 0;     static unsigned *connectivity = 0;     static unsigned connectivity_size = 0;     static int current_mode = 0; // 0: None, 1: Focused-Locality, 2: Diffuse-Global     static uint64_t double_until = 0;      // Memory management for connectivity tracking     if (!connectivity || VARS > connectivity_size) {         unsigned new_size = VARS + 1000;         unsigned *new_conn = (unsigned *)kissat_calloc(solver, new_size, sizeof(unsigned));         if (connectivity) {             kissat_free(solver, connectivity, connectivity_size * sizeof(unsigned));         }         connectivity = new_conn;         connectivity_size = new_size;     }      // Update connectivity map every 256 conflicts     static uint64_t last_conn_update = 0;     if (CONFLICTS >= last_conn_update + 256) {         for (unsigned i = 0; i < VARS; i++) connectivity[i] = 0;         for (all_clauses(c)) {             if (c->redundant && !c->garbage) {                 for (all_literals_in_clause(lit, c)) {                     connectivity[IDX(lit)]++;                 }             }         }         last_conn_update = CONFLICTS;     }      // Step 1: CLI Calculation every 1024 conflicts     if (CONFLICTS >= last_cli_checkpoint + 1024) {         double total_cli = 0;         unsigned red_count = 0;         for (all_clauses(c)) {             if (c->redundant && !c->garbage) {                 double avg_idx = 0;                 for (all_literals_in_clause(lit, c)) avg_idx += (double)IDX(lit);                 avg_idx /= c->size;                  double clause_dist = 0;                 for (all_literals_in_clause(lit, c)) {                     double d = (double)IDX(lit) - avg_idx;                     clause_dist += (d < 0 ? -d : d);                 }                 total_cli += (clause_dist / c->size) / (VARS > 0 ? VARS : 1);                 red_count++;             }         }         double cli = (red_count > 0) ? (total_cli / red_count) : 0;          // Step 2: Focused-Locality Mode (CLI < 0.25)         if (cli < 0.25) {             current_mode = 1;             kissat_set_option(solver, "decay", 750);             kissat_set_option(solver, "restart", 1);             kissat_set_option(solver, "restartint", 20);             kissat_set_option(solver, "stable", 1);         }         // Step 3: Diffuse-Global Mode (CLI > 0.6)         else if (cli > 0.6) {             current_mode = 2;             kissat_set_option(solver, "decay", 980);             kissat_set_option(solver, "stable", 0);             kissat_set_option(solver, "target", 2);         }         last_cli_checkpoint = CONFLICTS;     }      // Step 4: LBD Slope & Target-Best Synchronization every 2048 conflicts     if (CONFLICTS >= last_lbd_checkpoint + 2048) {         double current_avg_lbd = 0;         unsigned red_count = 0;         for (all_clauses(c)) {             if (c->redundant && !c->garbage) {                 current_avg_lbd += c->glue;                 red_count++;             }         }         current_avg_lbd /= (red_count > 0 ? red_count : 1);         double slope = current_avg_lbd - last_avg_lbd;          if (current_mode == 1 && slope > 0) {             unsigned threshold = (red_count / 10);             for (all_variables(i)) {                 if (ACTIVE(i) && connectivity[i] > threshold) {                     solver->phases.target[i] = solver->phases.best[i];                 }             }             if (double_until <= CONFLICTS) {                 solver->scinc *= 2.0;             }             double_until = CONFLICTS + 512;         }         last_avg_lbd = current_avg_lbd;         last_lbd_checkpoint = CONFLICTS;     }      // Manage VSIDS increment doubling timer     if (double_until > 0 && CONFLICTS >= double_until) {         solver->scinc /= 2.0;         double_until = 0;     } } \end{lstlisting} \end{codebox}  \subsection{Example 2: Kissat\_Corephase Solver, reset\_phases\_mab Function} \paragraph{Parent: Natural Language} Stratified Activity-Based Phase Re-anchoring: Step 1: Calculate two activity thresholds, T\_high (85th percentile) and T\_low (35th percentile), by performing a quickselect or sampling the current EVSIDS activity scores of all variables. Step 2: Iterate through every variable v in the solver. Step 3: If the activity of v is greater than T\_high, set its saved phase to the logical negation of its current target phase (!phases->target[v]). This forces the solver to explore the opposite side of the most contentious variables. Step 4: If the activity of v is between T\_low and T\_high (inclusive), set its saved phase to the 'best' phase (phases->best[v]), which represents the assignment held during the largest consistent trail found so far. Step 5: If the activity of v is below T\_low, set its saved phase to its current target phase (phases->target[v]) to maintain search stability in the less-frequently updated backbone of the formula.  \paragraph{Parent: Code Implementation} \begin{codebox} \begin{lstlisting}[style=cstyle] #include "autarky.h" #include "backtrack.h" #include "decide.h" #include "internal.h" #include "logging.h" #include "print.h" #include "rephase.h" #include "report.h" #include "terminate.h" #include "walk.h"  #include <inttypes.h> #include <string.h> #include <math.h>  static char reset_phases (kissat *solver) {   solver->rephased.count++;   const unsigned vars = VARS;   const value initial_phase = INITIAL_PHASE;    if (vars > 0) {     static double samples[100];     int n = 0;     unsigned step = vars / 100;     if (step == 0) step = 1;      // Step 1: Sample EVSIDS activity scores to calculate thresholds     for (unsigned i = 0; i < vars && n < 100; i += step) {       samples[n++] = solver->scores.score[i];     }      // Sort samples (simple selection sort for small n)     for (int i = 0; i < n - 1; i++) {       for (int j = i + 1; j < n; j++) {         if (samples[i] > samples[j]) {           double t = samples[i];           samples[i] = samples[j];           samples[j] = t;         }       }     }      // Calculate thresholds for 35th and 85th percentiles     double T_low  = (n > 0) ? samples[(35 * (n - 1)) / 100] : 0;     double T_high = (n > 0) ? samples[(85 * (n - 1)) / 100] : 0;      value *saved        = solver->phases.saved;     const value *best   = solver->phases.best;     const value *target = solver->phases.target;      // Steps 2-5: Stratified phase re-anchoring     for (unsigned v = 0; v < vars; v++) {       double act = solver->scores.score[v];       if (act > T_high) {         // Step 3: High activity -> negate current target phase         value t = target[v] ? target[v] : initial_phase;         saved[v] = -t;       } else if (act >= T_low) {         // Step 4: Medium activity -> set to best phase         saved[v] = best[v] ? best[v] : initial_phase;       } else {         // Step 5: Low activity -> maintain target phase         saved[v] = target[v] ? target[v] : initial_phase;       }     }   }    char type = 'B';   INC (rephased_best);  #ifndef QUIET   char const *type_as_string = 0;   switch (type) { #define REPHASE(NAME, TYPE, INDEX) \     case TYPE: type_as_string = #NAME; break;     REPHASES #undef REPHASE   }   assert (type_as_string);   kissat_phase (solver, "rephase", GET (rephased),     "%s phases in %s search mode (stratified activity-based)",     type_as_string, solver->stable ? "stable" : "focused"); #endif    solver->rephased.last = type;   solver->rephased.last_index = 0;   LOG ("copying saved phases as target phases");   memcpy (solver->phases.target, solver->phases.saved, vars);   UPDATE_CONFLICT_LIMIT (rephase, rephased, NLOGNLOGNLOGN, false);   kissat_reset_target_assigned (solver);   if (type == 'B')     kissat_reset_best_assigned (solver);   return type; } \end{lstlisting} \end{codebox}  \paragraph{Mutant 1: Natural Language} Agility-Sensitive Dynamic Threshold Re-anchoring: Step 1: Calculate dynamic activity thresholds T\_high and T\_low by adjusting the base 85th and 35th percentiles based on the solver's current agility (the exponential moving average of phase flips). Specifically, set the high-activity percentile to (100 - (15 * (1.0 + agility))) and the low-activity percentile to (35 * (1.0 - agility)), then extract these values from the current EVSIDS distribution using a partial quicksort on the activity heap. Step 2: Iterate through every variable v in the solver. Step 3: If the activity of v is greater than T\_high, set its saved phase to the logical negation of its current target phase (!phases->target[v]). This forces the solver to explore the opposite side of the most contentious variables. Step 4: If the activity of v is between T\_low and T\_high (inclusive), set its saved phase to the 'best' phase (phases->best[v]), which represents the assignment held during the largest consistent trail found so far. Step 5: If the activity of v is below T\_low, set its saved phase to its current target phase (phases->target[v]) to maintain search stability in the less-frequently updated backbone of the formula.  \paragraph{Mutant 1: Code Implementation} \begin{codebox} \begin{lstlisting}[style=cstyle] #include "autarky.h" #include "backtrack.h" #include "decide.h" #include "internal.h" #include "logging.h" #include "print.h" #include "rephase.h" #include "report.h" #include "terminate.h" #include "walk.h"  #include <inttypes.h> #include <string.h> #include <math.h>  static char reset_phases (kissat * solver) {   const unsigned n = VARS;    // Step 1: Calculate dynamic thresholds via agility   double agility_val  = solver->agility;   double p_high_pct   = 100.0 - (15.0 * (1.0 + agility_val));   double p_low_pct    = 35.0  * (1.0 - agility_val);    double t_high = 0;   double t_low  = 0;    static unsigned ids[2097152];    if (n > 0 && n <= 2097152) {     for (unsigned i = 0; i < n; i++) ids[i] = i;      unsigned k_low  = (unsigned)((p_low_pct  / 100.0) * (n - 1));     unsigned k_high = (unsigned)((p_high_pct / 100.0) * (n - 1));      // Iterative Quickselect for T_low     int l = 0, r = (int)n - 1;     while (l < r) {       double pivot = solver->scores.scores[ids[(l + r) / 2]];       int i = l, j = r;       while (i <= j) {         while (solver->scores.scores[ids[i]] < pivot) i++;         while (solver->scores.scores[ids[j]] > pivot) j--;         if (i <= j) {           unsigned tmp = ids[i]; ids[i] = ids[j]; ids[j] = tmp;           i++; j--;         }       }       if ((int)k_low <= j) r = j;       else if ((int)k_low >= i) l = i;       else break;     }     t_low = solver->scores.scores[ids[k_low]];      // Iterative Quickselect for T_high     l = 0; r = (int)n - 1;     while (l < r) {       double pivot = solver->scores.scores[ids[(l + r) / 2]];       int i = l, j = r;       while (i <= j) {         while (solver->scores.scores[ids[i]] < pivot) i++;         while (solver->scores.scores[ids[j]] > pivot) j--;         if (i <= j) {           unsigned tmp = ids[i]; ids[i] = ids[j]; ids[j] = tmp;           i++; j--;         }       }       if ((int)k_high <= j) r = j;       else if ((int)k_high >= i) l = i;       else break;     }     t_high = solver->scores.scores[ids[k_high]];   }    // Steps 2-5: Iterate variables and re-anchor phases   for (unsigned v = 0; v < n; v++) {     double activity = solver->scores.scores[v];     if (activity > t_high) {       // Step 3: High activity -> negation of target phase       solver->phases.saved[v] = -solver->phases.target[v];     } else if (activity >= t_low) {       // Step 4: Medium activity -> best phase       value b = solver->phases.best[v];       if (b) solver->phases.saved[v] = b;     } else {       // Step 5: Low activity -> maintain target phase       solver->phases.saved[v] = solver->phases.target[v];     }   }    char type = 'B';   INC (rephased_best);  #ifndef QUIET   char const *type_as_string = 0;   switch (type) { #define REPHASE(NAME, TYPE, INDEX) \     case TYPE: type_as_string = #NAME; break;     REPHASES #undef REPHASE   }   assert (type_as_string);   kissat_phase (solver, "rephase", solver->statistics.rephased,     "%s phases in %s search mode (agility-sensitive thresholds)",     type_as_string, solver->stable ? "stable" : "focused"); #endif    solver->rephased.last = type;   memcpy (solver->phases.target, solver->phases.saved, VARS);   UPDATE_CONFLICT_LIMIT (rephase, rephased, NLOGNLOGNLOGN, false);   kissat_reset_target_assigned (solver);   if (type == 'B')     kissat_reset_best_assigned (solver);   return type; } \end{lstlisting} \end{codebox} \paragraph{Mutant 2: Natural Language} Mode-Adaptive Mid-Tier Phase Diversification: Step 1: Calculate two activity thresholds, T\_high (85th percentile) and T\_low (35th percentile), by performing a quickselect or sampling the current EVSIDS activity scores of all variables. Step 2: Iterate through every variable v in the solver. Step 3: If the activity of v is greater than T\_high, set its saved phase to the logical negation of its current target phase (!phases->target[v]). This forces the solver to explore the opposite side of the most contentious variables. Step 4: If the activity of v is between T\_low and T\_high (inclusive), set its saved phase to the 'best' phase (phases->best[v]) if the solver is currently in 'stable' mode; however, if the solver is in 'focused' mode, set it to the negation of the best phase (!phases->best[v]) to aggressively diversify the search in the mid-activity range. Step 5: If the activity of v is below T\_low, set its saved phase to its current target phase (phases->target[v]) to maintain search stability in the less-frequently updated backbone of the formula.  \paragraph{Mutant 2: Code Implementation} \begin{codebox} \begin{lstlisting}[style=cstyle] #include "autarky.h" #include "backtrack.h" #include "decide.h" #include "internal.h" #include "logging.h" #include "print.h" #include "rephase.h" #include "report.h" #include "terminate.h" #include "walk.h"  #include <inttypes.h> #include <string.h> #include <math.h>  static char reset_phases(kissat *solver) {   solver->rephased.count++;   const unsigned vars   = VARS;   const value initial   = INITIAL_PHASE;   const bool stable     = solver->stable;   const double *const scores = solver->scores.score;    // Step 1: Sample EVSIDS scores and sort to find percentile thresholds   double samples[100];   unsigned n = 0;   if (vars > 0) {     for (unsigned i = 0; i < 100; i++) {       unsigned idx = (unsigned)(((unsigned long long)i * vars) / 100);       if (idx >= vars) idx = vars - 1;       samples[n++] = scores[idx];     }     for (unsigned i = 0; i < n; i++) {       for (unsigned j = i + 1; j < n; j++) {         if (samples[i] > samples[j]) {           double tmp = samples[i];           samples[i] = samples[j];           samples[j] = tmp;         }       }     }   }    double T_low  = (n > 0) ? samples[(35 * n) / 100] : 0;   double T_high = (n > 0) ? samples[(85 * n) / 100] : 0;    value *saved        = solver->phases.saved;   const value *best   = solver->phases.best;   const value *target = solver->phases.target;    // Steps 2-5: Mode-adaptive stratified rephase   for (unsigned v = 0; v < vars; v++) {     double act = scores[v];     value t = target[v] ? target[v] : initial;     value b = best[v]   ? best[v]   : initial;      if (act > T_high) {       // Step 3: High activity -> negate target phase       saved[v] = -t;     } else if (act >= T_low) {       // Step 4: Mid-tier -> stable mode uses best; focused mode negates best       saved[v] = stable ? b : -b;     } else {       // Step 5: Low activity -> maintain target phase       saved[v] = t;     }   }    const char type = 'B';  #ifndef QUIET   char const *type_as_string = 0;   switch (type) { #define REPHASE(NAME, TYPE, INDEX) \     case TYPE: type_as_string = #NAME; break;     REPHASES #undef REPHASE   }   assert(type_as_string);   kissat_phase(solver, "rephase", GET(rephased),                "%s phases in %s search mode",                type_as_string, solver->stable ? "stable" : "focused"); #endif    kissat_extremely_verbose(solver, "remembering last rephase type '%c' at "                                    "%s conflicts",                            type, FORMAT_COUNT(CONFLICTS));   solver->rephased.last = type;   solver->rephased.last_index = 0;   LOG("copying saved phases as target phases");   memcpy(solver->phases.target, solver->phases.saved, vars);   UPDATE_CONFLICT_LIMIT(rephase, rephased, NLOGNLOGNLOGN, false);   kissat_reset_target_assigned(solver);   if (type == 'B')     kissat_reset_best_assigned(solver);   return type; } \end{lstlisting} \end{codebox}  % \section{Example of Memory Bank}  % \section{Technical appendices and supplementary material} % Technical appendices with additional results, figures, graphs, and proofs may be submitted with the paper submission before the full submission deadline (see above). You can upload a ZIP file for videos or code, but do not upload a separate PDF file for the appendix. There is no page limit for the technical appendices.   % Note: Think of the appendix as ``optional reading'' for reviewers. The paper must be able to stand alone without the appendix; for example, adding critical experiments that support the main claims to an appendix is inappropriate.   %%%%%%%%%%%%%%%%%%%%%%%%%%%%%%%%%%%%%%%%%%%%%%%%%%%%%%%%%%%%  \newpage \input{checklist.tex} \end{document}
